% Standalone English manuscript. Compile twice with pdflatex.
% Full integration of the long-direction fourth-derivative refinement.
% All 51 affine certificates and the fourth/fifth derivative identities
% are included explicitly; no external TeX inputs are required.
% Status: candidate proof awaiting independent mathematical verification.
\documentclass[11pt,a4paper]{article}
\usepackage[OT1]{fontenc}
\usepackage[utf8]{inputenc}
\usepackage[margin=28mm]{geometry}
\usepackage{amsmath,amssymb,amsthm,mathtools}
\usepackage{booktabs,longtable,array}
\usepackage{microtype}
\usepackage[hidelinks]{hyperref}
\allowdisplaybreaks[2]
\numberwithin{equation}{section}
\newtheorem{theorem}{Theorem}[section]
\newtheorem{lemma}[theorem]{Lemma}
\newtheorem{proposition}[theorem]{Proposition}
\theoremstyle{remark}
\newtheorem{remark}[theorem]{Remark}
\newcommand{\N}{\mathbb N}
\newcommand{\Z}{\mathbb Z}
\newcommand{\R}{\mathbb R}
\newcommand{\Q}{\mathbb Q}
\newcommand{\eps}{\varepsilon}
\newcommand{\dd}{\,\mathrm d}
\newcommand{\norm}[1]{\left\lVert #1\right\rVert}
\DeclareMathOperator*{\Res}{Res}
\DeclareMathOperator{\Hess}{Hess}
\DeclareMathOperator{\Var}{Var}
\DeclareMathOperator{\supp}{supp}
\DeclareMathOperator{\grad}{grad}
\newcommand{\Leg}{\mathcal L}
\newcommand{\Rfine}{\mathcal R_{\mathrm{fine}}}
\newcommand{\pint}{p_{\mathrm{int}}}
\newcommand{\hcap}{h_{\mathrm{cap}}}
\newcommand{\hthreshold}{h_{\mathrm t}}
\newcommand{\Fzero}{F_0}
\newcommand{\Qone}{\mathcal Q_1}
\setlength{\emergencystretch}{2em}

\title{A candidate bound for the ternary divisor problem\\
\large A long-direction fourth-derivative refinement}
\author{}
\date{9 September 2026}

\begin{document}
\maketitle
\begin{abstract}
We give a detailed argument intended to establish the pointwise estimate
$\Delta_3(x)\ll_\eps x^{541/1209+\eps}$ for the classical symmetric ternary
divisor problem. The argument combines a truncated Voronoi formula,
differencing, successive one-dimensional van der Corput transforms on sharp
domains, and a fourth-derivative estimate along the longer terminal dual
direction. A sublevel decomposition treats zeros of that derivative by
using a uniform same-direction fourth/fifth derivative condition.
All stationary weights,
short dual ranges and integer choices of differencing parameters are retained.
Explicit polynomial certificates verify the needed model nondegeneracy, and
51 rational affine identities verify the global parameter inequalities.
\end{abstract}

\noindent\textbf{Working manuscript.}
This version integrates the long-direction refinement into the full argument,
including the boundary subregion, exact derivative certificates and all global
parameter certificates. This revision gives explicit uniform inverse charts,
sharp-fiber parameterizations, endpoint accounting, and a separate complexity
transfer lemma for the complete exact phase construction.

\tableofcontents

\section{The proposed result, conventions and external inputs}
\label{sec:inputs}
For $n\in\N$, put
\[
 d_3(n)=\#\{(a,b,c)\in\N^3:abc=n\},\qquad
 \Delta_3(x)=\sum_{n\le x}d_3(n)
   -\Res_{s=1}\frac{\zeta(s)^3x^s}{s}.
\]
Throughout, $e(t)=\exp(2\pi i t)$ and
\begin{equation}\label{eq:parameters}
 \theta=\frac{541}{1209},\qquad
 \rho=\frac{265}{403},\qquad
 u_0=\frac{138}{403}=3\theta-1.
\end{equation}
These constants satisfy
\begin{equation}\label{eq:endpoint-identities}
 \frac23-\frac\rho3=\theta,
 \qquad \frac13+\frac{u_0}{3}=\theta.
\end{equation}

\begin{theorem}[Proposed bound]\label{thm:main}
For every $\eps>0$ and all sufficiently large real $x$,
\[
 \Delta_3(x)\ll_\eps x^{541/1209+\eps}.
\]
\end{theorem}

The proof occupies Sections~\ref{sec:operations}--\ref{sec:assembly} and the
two appendices. No unproved spacing or moment estimate is an input.

\subsection{Uniformity and scales}
All normalized phase variables lie in fixed compact boxes with slightly larger
open neighborhoods. When positive variables are used, these neighborhoods stay
away from zero. Domains are semialgebraic sets of bounded description
complexity: the number and degrees of their defining polynomial conditions are
bounded independently of the large parameters. Their real coefficients and
sharp endpoints may vary. A fixed number of charts, Boolean operations,
projections and rectangular cuts is allowed.

The constants in $\ll$ may depend on these fixed complexity and derivative
bounds, and on a small positive auxiliary parameter $\eta$, but not on $x$.
All physical scales have at most polynomial size in $x$. Fixed powers of
$\log x$ will be absorbed into $x^\eta$, with $\eta$ reduced finitely often
before the final choice of $\eps$. In intermediate displays, the same symbol
$x^\eta$ may denote such a rescaled loss.

Dyadic lengths and bounding rectangles are understood up to fixed factors.
Whenever a one-dimensional theorem requires its length parameter to be at
least the actual interval length, take a sufficiently large fixed multiple of
the nominal ambient scale. This changes only the comparison constants. In
particular, the length of an individual sharp fiber need not be comparable to
the ambient scale.

In the global argument we distinguish the physical phase scale from its exact
exponent budget:
\begin{equation}\label{eq:F-convention}
 R=27x,\qquad F=(RN)^{1/3}=3\Fzero,
 \qquad \Fzero=(xN)^{1/3}.
\end{equation}
Exact stationary formulas retain $F$. Nominal parameter choices use $\Fzero$.
Also, $\Pi$ denotes a nominal product of shift lengths, while
$\pint=q_1q_2$ denotes its integer realization. Only comparability, not equality,
will be used between these two products.

\subsection{Analytic inputs}
We use the following published estimates, in the indicated normalizations.

First, Lemma~3 of Ivi\'c--Zhai~\cite{IZ}, specialized to $k=3$, gives
\begin{equation}\label{eq:voronoi}
 \Delta_3(x)=\frac{x^{1/3}}{\pi\sqrt3}
 \sum_{n\le Y}d_3(n)n^{-2/3}
       \cos\bigl(6\pi(xn)^{1/3}\bigr)
 +O_\eta\bigl(x^{2/3+\eta}Y^{-1/3}\bigr),
 \qquad 1\ll Y\ll x.
\end{equation}
An endpoint convention at an integer $x$ changes the left side by
$O(d_3(x))=O_\eta(x^\eta)$, which is harmless here.

Second, the one-dimensional B transform in Vandehey~\cite[Theorem~1.1]{Vand}
applies on an interval $[a,b]$ with ambient length $M\ge b-a$ if
\[
 f''\asymp T/M^2>0,\qquad
 |f'''|\ll T/M^3,\qquad |f''''|\ll T/M^4.
\]
For a real amplitude $g$ of bounded variation, it replaces the sum by the
stationary sum with terms
\[
 \frac{g(x_r)}{\sqrt{f''(x_r)}}
 e\bigl(f(x_r)-rx_r+1/8\bigr),\qquad f'(x_r)=r,
\]
and an error
\begin{equation}\label{eq:one-B-error}
 \ll \bigl(\Var(g)+|g(a)|\bigr)
 \left(\frac{M}{\sqrt T}
 +\log\bigl(2+f'(b)-f'(a)\bigr)\right).
\end{equation}
Negative curvature is handled by conjugation. In unweighted applications,
changes between open, closed and half-weight endpoint conventions cost
$O(1+M/\sqrt T)$. We explicitly retain this cost. The statement permits an
ambient $M$ larger than the fiber length.

Third, the classical third-derivative estimate, as recorded in
Heath-Brown~\cite[equation~(1), $k=3$]{HB}, is
\begin{equation}\label{eq:third-test}
 \sum_{n\in I}e(f(n))
 \ll M\lambda^{1/6}+M^{1/2}\lambda^{-1/6}
 \quad\text{if}\quad |f'''|\asymp\lambda,
 \quad |I|\le M,\quad M\ge1.
\end{equation}
The derivative has a fixed sign on each interval under consideration. For
shorter fibers one may first use their actual length and then enlarge the
bound to $M$; isolated points are handled separately when needed.

We also use the fourth-derivative instance of the same classical estimate,
Heath-Brown~\cite[equation~(1), $k=4$]{HB}:
\begin{equation}\label{eq:fourth-test}
 \sum_{n\in I}e(f(n))
 \ll M\lambda^{1/14}+M^{3/4}\lambda^{-1/14},
 \qquad |f^{(4)}|\asymp\lambda,\quad |I|\le M,\quad M\ge1.
\end{equation}
The same endpoint and sign conventions apply. If $\lambda$ is bounded below
by a positive constant, the first term already covers the trivial estimate.

Finally, we use real quantifier elimination with a format bound depending
only on the input format; see Basu~\cite[Section~2]{Basu}.
Lemmas~\ref{lem:sa-calculus} and~\ref{lem:sharp-fibers} below spell out
the resulting uniform graph, projection and fiber assertions. These
facts do not supply an estimate for lattice points near a curved boundary.

\section{Elementary operations and a sharp two-dimensional B inequality}
\label{sec:operations}

\begin{lemma}[A inequality]\label{lem:A}
Let $D$ be a finite lattice set in a rectangle of side lengths $O(M_1),O(M_2)$,
where $M_1,M_2\ge1$, and let $1\le q_j\ll M_j$ be integers. Then
\begin{align}\label{eq:A2-general}
 \left|\sum_{z\in D}e(f(z))\right|^2
 &\ll \frac{(M_1M_2)^2}{q_1q_2}\notag\\
 &\quad+\frac{M_1M_2}{q_1q_2}
 \sum_{\substack{|r_j|<q_j\\r\ne0}}
 \left|\sum_{z\in D\cap(D-r)}e\bigl(f(z+r)-f(z)\bigr)\right|.
\end{align}
The triangular weights may instead be retained before absolute values.
\end{lemma}
\begin{proof}
Extend the original summand by zero. Average over the $q_1q_2$ shifts in a
rectangle and apply Cauchy's inequality on the common enlarged support, which
has $O(M_1M_2)$ points. Expanding the square gives multiplicities
$(q_1-|r_1|)(q_2-|r_2|)$. The diagonal gives the first term; dividing the
off-diagonal multiplicities by $q_1q_2$ gives weights at most one. No
regularity of the boundary is required. If $q_1=1$, the first shift coordinate
is identically zero.
\end{proof}

\begin{lemma}[Rectangular Abel summation]\label{lem:Abel}
Let $a$ be $C^2$ on an ambient rectangle with arbitrary positive side lengths
$M_1,M_2$. For any finite base array supported on $D$,
\begin{equation}\label{eq:Abel}
 \left|\sum_{z\in D}a(z)e(f(z))\right|
 \ll V_{\rm rect}(a)
 \sup_J\left|\sum_{z\in D\cap J}e(f(z))\right|,
\end{equation}
where $J$ ranges over rectangles and
\[
 V_{\rm rect}(a)=\norm a_\infty
 +M_1\norm{a_1}_\infty+M_2\norm{a_2}_\infty
 +M_1M_2\norm{a_{12}}_\infty.
\]
Fixed coefficients may be retained in the base array on both sides.
\end{lemma}
\begin{proof}
Apply integration by parts in each coordinate to the finite measure defined
by the base array. The boundary, edge and interior terms have exactly the
four variation costs displayed above. Keep $1_D$ inside every partial sum.
The proof works when a side length is less than one; no extension to a
rectangle of side length one is necessary.
\end{proof}

\subsection{Uniform formulas, sharp fibers and inverse charts}

We use the following precise convention for complexity. A real formula has
\emph{format at most $B$} if it uses at most $B$ free and quantified scalar
variables, at most $B$ polynomial conditions and Boolean or quantifier
operations, and polynomials of degree at most $B$. Real coefficients are
unrestricted. Parameters, including cut endpoints, are free variables.
Descriptions with arbitrary coefficients fit this convention by treating
the coefficients as parameters; only finitely many Boolean patterns are
possible at a fixed format bound.
By quantifier elimination over the reals, a formula of format $B$ has an
equivalent quantifier-free formula whose format is at most a number $q(B)$.
Only the existence of $q(B)$ is used; see \cite[Section~2]{Basu}.
For a function, complexity means the format of its graph. Thus the
``algebraic families'' below are smooth functions with uniformly
semialgebraic graphs, with their real branches included in those graphs.

\begin{lemma}[Uniform semialgebraic calculus]\label{lem:sa-calculus}
Fix a format bound $B$, a derivative order $r$, and a number $L$ of
operations. Starting with sets and function graphs of format at most $B$,
each of the following operations, and any sequence of at most $L$ of them,
produces a graph or set of bounded format, depending only on $B,r,L$:
Boolean operations, coordinate projections, images and preimages under
the given maps, affine changes and translations, arithmetic operations
on their natural domains, derivatives of order at most $r$ where they
exist, and inverse maps restricted to a specified injective branch.
Restrictions by derivative inequalities are included. The same conclusion
holds for a difference quotient at any nonzero step, uniformly as that
step tends to zero, and for its derivative extension at zero.
Every one-dimensional fiber of any resulting set has a uniformly bounded
number of interval and point components.
\end{lemma}
\begin{proof}
Composition and images are expressed by existentially quantified graph
variables. Addition and multiplication add the equations $w=y_1+y_2$
and $w=y_1y_2$; division adds $wy_2=y_1$, $y_2\ne0$.
The graph of the inverse of $H$ on a specified set $C$ is precisely
\[
 \{(y,z):z\in C,\ (z,y)\in\operatorname{graph}(H)\}.
\]
Injectivity makes this relation single-valued. No choice of an unspecified
root is part of this operation.

For completeness, if $h$ is differentiable on an open spatial domain,
$v=\partial_i h(z)$ is expressed by the usual limit formula
\[
 \forall\epsilon>0\ \exists\delta>0\ \forall t\quad
 0<|t|<\delta\ \Longrightarrow\
 |h(z+te_i)-h(z)-tv|<\epsilon|t|,
\]
with $\delta$ also required to keep the indicated segment in the domain.
Replace function values by graph variables. The last inequality is the
polynomial inequality
$(y_t-y_0-tv)^2<\epsilon^2t^2$; the domain condition also has a fixed
formula. Iteration gives every derivative of the fixed order $r$.
For $a\ne0$, the difference quotient has graph
\begin{equation}\label{eq:sa-quotient-graph}
 \exists y_0,y_1:\quad
 (z,y_0)\in\operatorname{graph}(h),\quad
 (z+ae_i,y_1)\in\operatorname{graph}(h),\quad
 aw=y_1-y_0.
\end{equation}
At $a=0$ use the derivative graph as a separate parameter piece. A general
specified direction is handled by replacing $e_i$ by that direction.

Every construction adds a bounded number of variables, conditions and
degrees. Apply $q$ after each of the finitely many operations. This gives
a recursive bound depending only on $B,r,L$, without bounds on coefficient
sizes or on nonzero denominators. After fixing all but one spatial
coordinate, a quantifier-free description involves at most $s$ univariate
polynomials of degrees at most $d$, for uniformly bounded $s,d$.
Discard the identically zero polynomials. Their at most $sd$ real roots
cut the line into at most $2sd+1$ points and open intervals on which every
condition is constant. This proves the fiber assertion, including
degenerate parameter values.
\end{proof}

\begin{lemma}[Sharp-fiber parameterization and endpoint accounting]
\label{lem:sharp-fibers}
Let $S_\alpha\subset\R^2$ be a bounded semialgebraic family of fixed
format; $\alpha$ includes all cut endpoints. For $i=1,2$, define its
coordinate-fiber interior by
\begin{equation}\label{eq:fiber-interior}
 \mathfrak I_iS_\alpha
 =\{z:\exists\delta>0\ \forall t\ (|t|<\delta
                \Longrightarrow z+te_i\in S_\alpha)\}.
\end{equation}
This is a family of bounded format contained in $S_\alpha$.
There is a uniform integer $N_*$ such that every $i$-fiber of
$\mathfrak I_iS_\alpha$ is a disjoint union of at most $N_*$ open
intervals, and every fiber of $S_\alpha\setminus\mathfrak I_iS_\alpha$
contains at most $N_*$ points.

The interval endpoints can be labeled by a bounded number of
semialgebraic functions of $\alpha$ and the other coordinate, on
semialgebraic parameter pieces of bounded format. Consequently, if the
other coordinate ranges over an interval of length $O(M)$, replacing
$S_\alpha$ by $\mathfrak I_iS_\alpha$ in a lattice sum with summands of
absolute value at most $A$ costs $O(A(1+M))$.
\end{lemma}
\begin{proof}
Formula \eqref{eq:fiber-interior} and Lemma~\ref{lem:sa-calculus} give
bounded format. Apply the root count in that lemma to a fiber. The
interior consists of its open interval components. What remains consists
only of included interval endpoints and isolated points, with a uniform
bound on their total number. This includes fibers of dimension zero and
empty fibers.

Here is an explicit endpoint parameterization. Fiberwise closure is
defined by
\[
 z\in\operatorname{cl}_i S_\alpha
 \quad\Longleftrightarrow\quad
 \forall\epsilon>0\ \exists t\quad
 |t|<\epsilon,\quad z+te_i\in S_\alpha.
\]
The fiberwise boundary
$\operatorname{cl}_i S_\alpha\setminus\mathfrak I_iS_\alpha$
is a finite set of uniformly bounded cardinality. Partition the parameter
space, including the other coordinate, by that cardinality $k$. Its
ordered elements $b_1<\cdots<b_k$ have graphs defined by membership in
this boundary and the assertion that exactly $j-1$ boundary elements lie
below $b_j$. The latter assertion uses at most $N_*$ quantified variables:
list these elements and require every smaller boundary element to be one
of them. Thus all these graphs have bounded format.

For adjacent boundary elements, membership of their midpoint decides
whether the entire open gap belongs to the fiber. Membership of each
$b_j$ decides its endpoint convention. Partition further by these
finitely many flags. Because the fibers are bounded, the two exterior
gaps are empty. This describes every fiber by ordered open intervals and
individually included endpoints. The labels may be discontinuous; the
argument neither differentiates them nor assumes continuity in any
parameter. Summing the uniform endpoint count over the $O(1+M)$ integer
values of the other coordinate proves the final assertion.
\end{proof}

\begin{lemma}[Common buffered inverse branches]\label{lem:buffered-inverses}
Let $K\subset\R^2$ be a fixed compact rectangle, and let $\Omega$ contain
a fixed neighborhood of $K$. Suppose $F_0,F\in C^m(\Omega)$, $m\ge4$,
have norms at most $M$ and
\[
 |\det\Hess F_0|\ge c,\qquad
 \max(|F_{0,11}|,|F_{0,22}|)\ge c
\]
on that neighborhood. There are positive constants $r,R,\delta_0,b$
and an integer $J$, depending only on $K,\Omega,M,c$, with the following
properties whenever $\norm{F-F_0}_{C^2}\le\delta_0$.

There are $J$ source rectangles $C_j$ covering $K$, each contained with
a fixed buffer in a ball $B_j=B(z_j,R)$ with
$\overline{B_j}\subset\Omega$, and a pivot
$i_j\in\{1,2\}$ on each ball. Both $F$ and $F_0$ have nonzero Hessian
determinant and nonzero $i_j$-th pure second derivative throughout $B_j$,
with uniform lower bounds and the same respective signs. For $H=F,F_0$
write
\[
 M_H(z)=(H_i(z),z_{3-i}),\qquad N_H(z)=\grad H(z)
 \quad(i=i_j).
\]
Both maps are injective on $B_j$. There are axis-parallel rectangles
$A_j^-\Subset A_j^+$ and $Q_j^-\Subset Q_j^+$, common to $F,F_0$, such
that both inverse maps are defined on neighborhoods of $\overline{A_j^+}$
and $\overline{Q_j^+}$, respectively, and
\begin{align}
 M_H(C_j)&\subset A_j^-,&N_H(C_j)&\subset Q_j^-,\notag\\
 N_H^{-1}(Q_j^+)&\subset M_H^{-1}(A_j^+)\subset B_j.
 \label{eq:compatible-rectangles}
\end{align}
All the rectangles have fixed positive side lengths, and each inner-to-outer
buffer is at least $b$. The inverse maps have uniform $C^{m-1}$ bounds.
The two Legendre functions, formed using these specified inverse branches,
satisfy, for $1\le k\le m-2$,
\begin{equation}\label{eq:inverse-stability}
 \norm{N_F^{-1}-N_{F_0}^{-1}}_{C^k(Q_j^+)}
 +\norm{\Leg F-\Leg F_0}_{C^k(Q_j^+)}
 \le C_k\norm{F-F_0}_{C^{k+1}(\Omega)}.
\end{equation}
The analogous inverse comparison holds on $A_j^+$.

If the phase graphs form a family of bounded format, all the chart
choices, inverse graphs and rectangles above can be parameterized with
bounded format. The source centers can be fixed in advance. The pivot
may depend on the model parameter. No chart depends on a sharp domain,
on a cut endpoint, or on a physical aspect ratio.
\end{lemma}
\begin{proof}
Increase $M\ge1$ to bound the operator norms of all Jacobians and their
first derivatives used here. At a source center, the least singular value
of $\Hess F_0$ is at least $c/M$. Choose a pivot with
$|F_{0,ii}(z_j)|\ge c$. The absolute determinant of the mixed Jacobian is
$|F_{0,ii}|$; its norm is bounded, so it too has a uniform least singular
value $\sigma>0$.
By decreasing a fixed $R>0$, arrange that $B(z_j,R)\subset\Omega$,
the selected pure derivative is at least $c/2$ in absolute value there,
and both model Jacobians vary from their central values by at most
$\sigma/8$. Decrease $\delta_0$ so that the perturbed Jacobians differ
from the model Jacobians by at most $\sigma/8$. These choices are uniform
in the center and parameter. They also preserve determinant and pivot
signs throughout the connected ball.

For either map $H_\nu=M_{F_\nu}$ or $N_{F_\nu}$, where $\nu=0,1$
denotes model and exact, put $A=DH_0(z_j)$. Then
\[
 \norm{DH_\nu(z)-A}\le\sigma/4,\qquad \norm{A^{-1}}\le\sigma^{-1}
 \quad(z\in B_j).
\]
Integration on the segment between $z,w\in B_j$ gives
\begin{equation}\label{eq:inverse-co-lipschitz}
 |H_\nu(z)-H_\nu(w)|\ge(3\sigma/4)|z-w|.
\end{equation}
In particular both maps are injective on the whole ball. If
$|y-H_0(z_j)|\le\sigma R/8$ and
$|H_1(z_j)-H_0(z_j)|\le\sigma R/8$, the map
\[
 z\longmapsto z+A^{-1}(y-H_\nu(z))
\]
is a contraction of the closed radius-$R$ ball into itself: its Lipschitz
constant is at most $1/4$, and its displacement of the center is at most
$R/4$. Its unique fixed point lies within $R/3$ of $z_j$.
Thus the two images contain the same target ball of radius $\sigma R/8$,
and their inverses there take values strictly inside $B_j$.

We give the nesting choices to ensure compatibility of the two inverses.
First choose concentric mixed rectangles $A^-\Subset A^+$ about
$M_{F_0}(z_j)$ whose closures lie strictly inside that common mixed
target ball. If $A^+$ has half-side $a$, choose the half-side $q$ of
$Q^+$ about $N_{F_0}(z_j)$ sufficiently small in terms of $a,M,\sigma$.
Indeed \eqref{eq:inverse-co-lipschitz} gives
\[
 |N_H^{-1}(y)-z_j|
 \le \frac4{3\sigma}
       \bigl(|y-N_{F_0}(z_j)|+C\delta_0\bigr).
\]
Consequently
\[
 |M_H(N_H^{-1}(y))-M_{F_0}(z_j)|
 \le \frac{4M}{3\sigma}(\sqrt2q+C\delta_0)+C\delta_0.
\]
Choose $q>0$ so that the first term with $\delta_0=0$ is less than
$a/4$, and then decrease $\delta_0$ so that the whole expression is less
than $a/2$. Choose $Q^+$ itself strictly within the common gradient
target ball, and take $Q^-$ and $A^-$ to have half the respective outer
half-sides. Finally choose a source half-side $r>0$ with
$\sqrt2Mr+C\delta_0<\min(a,q)/4$, decreasing $\delta_0$ again if needed.
Then $C_j=z_j+[-r,r]^2$ has the asserted image inclusions and fixed
buffers. The same estimates hold on slightly larger target rectangles,
so inverses are defined on neighborhoods of their closures. A finite
$r$-net of $K$, fixed before any parameter is chosen, provides the source
centers and a uniform number $J$.

Implicit differentiation of $H_\nu(I_\nu(y))=y$ first gives
$DI_\nu=(DH_\nu\circ I_\nu)^{-1}$. At order $k$, the term containing
$D^kI_\nu$ has coefficient $DH_\nu\circ I_\nu$; every other term is a
finite product of derivatives $D^aH_\nu$, $a\le k$, and lower derivatives
of $I_\nu$. This proves the uniform $C^{m-1}$ bounds by induction.
For comparison, \eqref{eq:inverse-co-lipschitz} first bounds
$|I_1-I_0|$ by $C\norm{H_1-H_0}_\infty$. Subtract the two differentiated
identities and induct. Evaluation of $D^kH_0$ at $I_1$ rather than $I_0$
is controlled by its bounded next derivative. For $k\le m-2$ this gives
$\norm{I_1-I_0}_{C^k}\le C_k\norm{F-F_0}_{C^{k+1}}$.
The formula
$\Leg H(y)=H(N_H^{-1}(y))-y\cdot N_H^{-1}(y)$ and
$\grad\Leg H=-N_H^{-1}$ give \eqref{eq:inverse-stability}.

Finally, use the first pivot when $|F_{0,11}(z_j)|\ge c$, and otherwise
the second. This is a fixed semialgebraic rule. Target centers are graph
values of the two model maps; all radii are fixed constants. An inverse
graph is specified by $z\in B_j$ and $y=M_H(z)$ or $y=N_H(z)$.
Lemma~\ref{lem:sa-calculus} gives the asserted format bounds, including
all parameter-dependent choices. The lower bounds proved above make
these graph relations unique branches.
\end{proof}


\subsection{The sharp B transformation}

For $U,V,P\ge1$, write
\begin{align}
 \mathcal R(U,V,P)
 &=\left(\frac{UV}{\sqrt P}+U+V+\sqrt P+\frac{UV}{P}\right)
       \log(2+U+V+P),\label{eq:R}\\
 \Rfine(U,V,P)
 &=\left(\frac{UV}{\sqrt P}+V+\sqrt P+\frac{UV}{P}\right)
       \log(2+U+V+P).\label{eq:Rfine}
\end{align}

\begin{lemma}[Sharp two-dimensional B inequality]\label{lem:sharpB}
Let $f(s,t)=P\mathcal F(s/U,t/V)$, with $U,V,P\ge1$.
On a fixed neighborhood of a compact normalized rectangle $K$, suppose
$\mathcal F$ is uniformly $C^{20}$ bounded, its graph has bounded format,
and
\[
 |\det\Hess\mathcal F|\ge c>0,\qquad
 \max(|\mathcal F_{11}|,|\mathcal F_{22}|)\ge c.
\]
Let $D\subset\{(Uz_1,Vz_2):z\in K\}$ be any semialgebraic sharp domain
of bounded format, with arbitrary real parameters and endpoint conventions.
Then
\begin{equation}\label{eq:sharpB}
 \left|\sum_{D\cap\Z^2}e(f)\right|
 \ll\frac{UV}{P}\sup_E\left|\sum_{E\cap\Z^2}e(f_j^*)\right|
       +\mathcal R(U,V,P).
\end{equation}
The supremum includes one of a bounded number of specified source charts
$j$ and two arbitrary real rectangles $J_1,J_2$. On that chart,
\begin{equation}\label{eq:sharp-domain-family}
 \begin{split}
 D_{1,j}&=\mathfrak I_1D^{(j)},\qquad
 A_{1,j}=T_{1,j}(D_{1,j})\cap J_1,\qquad
 A_{2,j}=\mathfrak I_2A_{1,j},\\
 E_{j,J_1,J_2}&=T_{2,j}(A_{2,j})\cap J_2.
 \end{split}
\end{equation}
Here this notation is for a first-coordinate pivot; the roles of the
coordinates are interchanged for the other pivot and the final frequency
labels are then restored. The maps are
$T_1(s,t)=(f_s(s,t),t)$ and $T_2(\xi,t)=(\xi,g_t(\xi,t))$, with $g$
the mixed phase on the specified inverse branch. The phase
$f_j^*(\xi,\zeta)=f(s,t)-\xi s-\zeta t$ uses the gradient inverse on
that same chart.

The entire family \eqref{eq:sharp-domain-family}, including the free
rectangle endpoints, has bounded format. Each $E$ is contained in the
gradient image of its source piece and in an inner inverse rectangle of
side lengths $O(P/U),O(P/V)$ with a fixed normalized buffer. If
$|\mathcal F_{11}|\ge c$ throughout, one can always use the order $s,t$
and replace $\mathcal R$ by $\Rfine$.
\end{lemma}
\begin{proof}
Apply Lemma~\ref{lem:buffered-inverses} with $F=F_0=\mathcal F$.
Write its normalized source rectangles as $C_1,\ldots,C_J$, and their
physical versions as $\mathcal C_j=\operatorname{diag}(U,V)C_j$.
Use the ordered partition
\begin{equation}\label{eq:chart-partition}
 D^{(j)}=(D\cap\mathcal C_j)\setminus\bigcup_{i<j}\mathcal C_i,
 \qquad D=\bigsqcup_{j=1}^J D^{(j)}.
\end{equation}
This assigns every real point, hence every lattice boundary point, exactly
once. Lemma~\ref{lem:sa-calculus} bounds the format of every piece.
Its closure stays inside the larger source ball where both inverse maps
and the curvature bounds are valid. We estimate each piece separately.

Fix a first-coordinate pivot and suppress $j$. The normalized mixed
coordinates are $(\xi U/P,t/V)$, while the normalized final coordinates
are $(\xi U/P,\zeta V/P)$. The mixed and gradient inverses exist on the
full physical rectangles obtained by scaling $A^+$ and $Q^+$ from
Lemma~\ref{lem:buffered-inverses}. Their branches agree wherever both
descriptions apply, by \eqref{eq:compatible-rectangles} and injectivity
on the source ball. In particular they agree on every set in
\eqref{eq:sharp-domain-family}.

Put $D_1=\mathfrak I_1D$. Lemma~\ref{lem:sharp-fibers} shows that
discarding $D\setminus D_1$ costs $O(1+V)$. For every integer $t$, the
remaining fiber is a union of at most $N_*$ disjoint open intervals
$(a,b)$. Its closure lies inside the buffered source environment and
\[
 |f_{ss}|\asymp P/U^2,\qquad
 |f_{sss}|\ll P/U^3,\qquad |f_{ssss}|\ll P/U^4.
\]
Apply the one-dimensional B theorem to $[a,b]$ with ambient length a
fixed multiple of $U$. Passing from closed or half-weight original
endpoints to $(a,b)$ costs at most two original summands. Passing to
the strictly interior stationary range $f_s((a,b),t)$ costs at most two
stationary summands, each of size $O(U/\sqrt P)$. Since $f_{ss}$ has
constant sign on the entire chart, the stationary point in this interval
is unique and all its Maslov factors equal $e(\operatorname{sgn}(f_{ss})/8)$.

Thus, with $L_0=\log(2+U+V+P)$, the exact first transformed expression is
\[
 \sum_{D\cap\Z^2}e(f)
 =e(\operatorname{sgn}(f_{ss})/8)\frac U{\sqrt P}
   \sum_{T_1(D_1)\cap\Z^2}a_1(\xi,t)e(g(\xi,t))+O(E_1),
\]
where
\begin{equation}\label{eq:E1}
 E_1\ll(1+V)(1+U/\sqrt P)L_0.
\end{equation}
This includes the B error, all original endpoint decisions, all dual
endpoint decisions, and all isolated original fibers. Disjoint source
intervals give disjoint open stationary intervals because $T_1$ is
injective. Hence the displayed dual sum has multiplicity one.

Writing $z=M_{\mathcal F}^{-1}(\xi U/P,t/V)$, the amplitude and mixed
phase are
\[
 a_1(\xi,t)=|\mathcal F_{11}(z)|^{-1/2},\qquad
 g(\xi,t)=f(s(\xi,t),t)-\xi s(\xi,t).
\]
The amplitude has bounded $C^2$ norm in normalized mixed coordinates on
the full outer rectangle. Lemma~\ref{lem:Abel} therefore bounds its
weighted sum by a constant times
\[
 \sup_{J_1}\left|\sum_{A_1\cap\Z^2}e(g)\right|,
 \qquad A_1=T_1(D_1)\cap J_1.
\]
The cuts may be taken within the mixed ambient rectangle; all their
endpoints are free real parameters. The variation bound remains uniform
when $P/U<1$, because the corresponding derivative factor is $U/P$ and
the ambient side factor is $P/U$.

On this same mixed environment, implicit differentiation gives
\begin{equation}\label{eq:mixed-curvature}
 g_{tt}=f_{tt}-\frac{f_{st}^2}{f_{ss}}
       =\frac{\det\Hess f}{f_{ss}},\qquad
 |g_{tt}|\asymp P/V^2.
\end{equation}
The third and fourth $t$ derivatives are $O(P/V^3)$ and $O(P/V^4)$:
this follows by differentiating the mixed inverse, whose normalized
derivatives are uniformly bounded by Lemma~\ref{lem:buffered-inverses}.
The sign of $g_{tt}$ is constant because both determinant and pivot signs
are constant throughout the source ball.

There are $O(1+P/U)$ integer values of $\xi$ in the mixed rectangle.
For each, replace the $t$ fiber of $A_1$ by its interior, setting
$A_2=\mathfrak I_2A_1$. Its discarded points cost $O(1)$ per $\xi$.
Apply the one-dimensional B theorem on each remaining open interval,
as above, now with ambient length a fixed multiple of $V$. Original
endpoint changes cost $O(1)$ and dual endpoint changes $O(V/\sqrt P)$
per interval. The number of intervals is uniformly bounded by
Lemma~\ref{lem:sharp-fibers}, uniformly in $J_1$. After multiplying by
the preceding Abel factor $O(U/\sqrt P)$, the total second error is
\begin{equation}\label{eq:E2}
 E_2\ll\frac U{\sqrt P}(1+P/U)(1+V/\sqrt P)L_0.
\end{equation}

Set $D_2=T_1^{-1}(A_2)$. The identity $g_t=f_t\circ T_1^{-1}$ gives
$T_2T_1=\grad f$ and therefore
\begin{equation}\label{eq:domain-provenance}
 E=T_2(A_2)\cap J_2=\grad f(D_2)\cap J_2
   \subset\grad f(D),\qquad D_2\subset D_1\subset D.
\end{equation}
The second stationary phase is exactly $f_j^*$ on the original gradient
branch. Its amplitude is $V/\sqrt P$ times the function
\[
 a_2(\xi,\zeta)
 =\left|\frac{\mathcal F_{11}}{\det\Hess\mathcal F}\right|^{1/2}
   \!\left(N_{\mathcal F}^{-1}(\xi U/P,\zeta V/P)\right).
\]
This function is defined and has bounded normalized $C^2$ norm on the
full outer gradient rectangle, before any sharp restriction. A second
application of Lemma~\ref{lem:Abel} gives the factor $O(UV/P)$ and the
second rectangular cut $J_2$. No endpoint half-weight remains as a
discontinuous amplitude: every such change was already charged in
$E_1$ or $E_2$.

Formula \eqref{eq:sharp-domain-family} is a parameterization of all the
sets over which we take a supremum. Each operation in it has bounded
format by Lemmas~\ref{lem:sa-calculus} and~\ref{lem:sharp-fibers}.
In particular the fiber interiors are taken as real sets for every real
value of the other coordinate, before restricting either coordinate to
integers. No set of integers, parameter-dependent enumeration of lattice
points, or union over all possible cuts is inserted into these formulas.
Lemma~\ref{lem:buffered-inverses} and
\eqref{eq:domain-provenance} place every $E$ in the designated inner
gradient rectangle, irrespective of both Abel cuts.

Finally, expansion of the two errors gives
\[
 E_1+E_2\ll
 \left(1+2V+\frac{2U}{\sqrt P}+\frac{UV}{\sqrt P}
             +\sqrt P+\frac{UV}{P}\right)L_0
 \ll\Rfine(U,V,P).
\]
Here $1\le V$ and $U/\sqrt P\le UV/\sqrt P$.
For the other pivot exchange $U,V$ in this calculation and restore
$(\xi,\zeta)=(f_s,f_t)$ at the end; the common remainder is
$\mathcal R$. Sum over the disjoint chart pieces in
\eqref{eq:chart-partition}. Their bounded number is absorbed in the
constant and their specified branches are included in the supremum.
This proves \eqref{eq:sharpB} with uniform constants for all sharp
parameters, endpoint coincidences, thin domains and short dual sides.
\end{proof}

\begin{remark}\label{rem:trivialdual}
The complete trivial dual main term in Lemma~\ref{lem:sharpB} is
\begin{equation}\label{eq:trivial-dual}
 \frac{UV}{P}(1+P/U)(1+P/V)=P+U+V+UV/P.
\end{equation}
The rectangular hull of a transformed domain is used for counting and
amplitude norms only; it does not replace the domain in an oscillating sum.
Any later A operation is applied after chart splitting, so that correlations
do not mix different inverse branches.
\end{remark}

\section{The model phase and exact finite differences}
\label{sec:model}
For a smooth function with an invertible local gradient, our Legendre
convention is
\[
 (\Leg G)(p,q)=G(u,v)-pu-qv,\qquad (p,q)=\grad G(u,v).
\]
In particular $\Hess(\Leg G)=-(\Hess G)^{-1}$.

\begin{lemma}[Uniform model nondegeneracy]\label{lem:model}
Let $u,v$ vary in a fixed positive compact box and let $(a,b)$ vary in a
compact set separated from $(0,0)$. Set
\begin{equation}\label{eq:G0}
 G_0(u,v)=\frac{a\sqrt u}{v}+\frac{bu^{3/2}}{v^2},
 \qquad \Phi_0=\Leg G_0,\quad K_0=\partial_1\Phi_0,
 \quad \Psi_0=\Leg K_0.
\end{equation}
On local inverse branches, the determinants of $\Hess G_0$ and $\Hess K_0$
are uniformly bounded away from zero. Each of these two Hessians has a
uniformly nonzero maximum of its pure second derivatives. Moreover,
\begin{equation}\label{eq:pure-third-max}
 \max\bigl(|\Psi_{0,111}|,|\Psi_{0,222}|\bigr)\ge c>0
\end{equation}
uniformly on the corresponding compact inverse environments. In addition,
the same second coordinate satisfies
\begin{equation}\label{eq:pure-four-five-model}
 \max\bigl(|\Psi_{0,2222}|,|\Psi_{0,22222}|\bigr)\ge c>0.
\end{equation}
\end{lemma}
\begin{proof}
Put
\begin{align*}
 D&=a^2v^2+4abuv+6b^2u^2=(av+2bu)^2+2b^2u^2,\\
 E&=a^2v^2+6abuv+18b^2u^2=(av+3bu)^2+9b^2u^2.
\end{align*}
Direct differentiation, with $\Hess K_0$ evaluated at $\grad G_0(u,v)$,
gives
\begin{equation}\label{eq:model-dets}
 \det\Hess G_0=-\frac{3D}{4uv^6},\qquad
 \det\Hess K_0=-\frac{160u^2v^{10}E}{27D^3}.
\end{equation}
These quantities are nonzero. The two pure second derivatives of $G_0$
could vanish simultaneously only if $av=3bu$ and $av=-3bu$, which is
impossible on the parameter set.

The pure second derivatives of $K_0$ have quartic homogeneous numerators,
and the pure third derivatives of $\Psi_0$ have degree-nine homogeneous
numerators. Appendix~\ref{app:model} gives the full polynomials, the
derivative identities producing them, and exact B\'ezout certificates
showing that each pair has no common direction zero. The direction $a=0$
is checked separately there. For \eqref{eq:pure-four-five-model},
the same appendix gives the degree-fourteen and degree-nineteen numerators
$Q_{14}$ and $Q_{19}$, their B\'ezout identity modulo $101$, and the
missing projective direction. This proves all the required noncoincidences.
Continuity and compactness give uniform positive lower bounds. All
normalized upper derivative bounds on buffered inverse charts follow
from the same determinant bounds and smoothness.
\end{proof}

\subsection{The exact family}
For small $e_0\ge0$, define
\begin{equation}\label{eq:Fe}
 \mathcal F_{e_0}(u,v)=\frac{u^{3/2}}v
       \left(\frac3{c^2+c+1}\right)^{3/2},
 \qquad c=(1+e_0u/v)^{1/3}.
\end{equation}
If $\omega=(\omega_1,\omega_2)$ is a unit vector and $\beta>0$ is small, put
\begin{equation}\label{eq:Gexact}
 G_{e_0,\beta,\omega}(u,v)
 =\frac{\mathcal F_{e_0}(u+\beta\omega_1,v+\beta\omega_2)
              -\mathcal F_{e_0}(u,v)}{\beta}.
\end{equation}
The continuous extension at $\beta=0$ is understood only for derivative
comparison. By the integral representation of a difference quotient,
\begin{equation}\label{eq:G-close}
 \norm{G_{e_0,\beta,\omega}-G_0}_{C^k}\ll e_0+\beta,
 \qquad a=\tfrac32\omega_1,\quad b=-\omega_2,
\end{equation}
for every fixed $k$, on a slightly enlarged positive compact set. Indeed
$\mathcal F_0=u^{3/2}/v$, and its directional derivative is precisely
\eqref{eq:G0}. The direction parameters form a fixed ellipse separated
from the origin. The constants in \eqref{eq:G-close} include the axis
directions. The radical family has uniformly bounded derivatives of every
fixed order in this neighborhood.

\begin{lemma}[Compatible inverse charts for the second difference]
\label{lem:transfer}
Let $G=G_{e_0,\beta,\omega}$, with $\omega$ a unit vector, and put
$\Phi=\Leg G$ using the branches specified below. There are a uniformly
bounded number of source charts and common exact/model inverse rectangles
$Q^-\Subset Q^\circ\Subset Q^+$, with uniform positive buffers, such
that on the common outer rectangles
\[
 \norm{\Phi-\Phi_0}_{C^{13}(Q^+)}\ll e_0+\beta.
\]
Every first-B domain, including both Abel cuts, belongs to its designated
$Q^-$ in normalized coordinates. There is a fixed $c_0>0$ such that for
$0<\lambda\le c_0$ the phase
\begin{equation}\label{eq:Phi-lambda}
 \Phi_\lambda(p,q)
 =\frac{\Phi(p+\lambda,q)-\Phi(p,q)}\lambda
 =\int_0^1\Phi_p(p+t\lambda,q)\dd t
\end{equation}
is defined on a neighborhood of $\overline{Q^\circ}$ and satisfies
\begin{equation}\label{eq:Phi-lambda-close}
 \norm{\Phi_\lambda-K_0}_{C^9(Q^\circ)}
 \ll e_0+\beta+\lambda,
 \qquad K_0=\Phi_{0,p}.
\end{equation}
Its Hessian determinant and pure second maximum are uniformly nonzero
for sufficiently small fixed upper bounds on $e_0,\beta,\lambda$.
There is a further uniformly bounded chart cover of $\overline{Q^-}$
with common buffered inverse rectangles for
$\grad\Phi_\lambda$ and $\grad K_0$. On each such rectangle, the specified
terminal Legendre functions obey
\[
 \norm{\Psi-\Psi_0}_{C^5}\ll e_0+\beta+\lambda,
 \qquad \Psi=\Leg\Phi_\lambda,\quad\Psi_0=\Leg K_0,
 \qquad \max(|\Psi_{2222}|,|\Psi_{22222}|)\ge c>0.
\]
On their respective buffered environments one also has
\begin{equation}\label{eq:uniform-high-regularity}
 \norm{G}_{C^{24}}+\norm{\Phi}_{C^{22}}
       +\norm{\Phi_\lambda}_{C^{21}}+\norm{\Psi}_{C^{20}}\le C.
\end{equation}
All constants and chart counts are independent of the small nonzero
steps and of $\omega$, including its axis directions. The sharp domain
and every cut endpoint are independent free parameters in these assertions.
\end{lemma}
\begin{proof}
Choose nested positive source rectangles in advance, with enough room
for the translations in \eqref{eq:Gexact}. The model parameters
$(a,b)=(3\omega_1/2,-\omega_2)$ range over a compact ellipse separated
from zero. Lemma~\ref{lem:model} supplies uniform determinant and pivot
margins on these rectangles. Also \eqref{eq:G-close} gives $C^{17}$
closeness of $G$ to $G_0$.

We first justify all upper regularity needed for the inverse estimates.
Formula \eqref{eq:Fe}, with its positive radical branches, has a uniformly
bounded $C^{25}$ norm on the larger positive rectangle. The identity
\[
 G(z)=\int_0^1\omega\cdot\grad\mathcal F_{e_0}
                    (z+t\beta\omega)\dd t
\]
gives a uniform $C^{24}$ bound for $G$. It also proves all the stated
comparison bounds without dividing an error by $\beta$. Decrease the
fixed upper bounds on $e_0,\beta$ to preserve the model Hessian margins.

Apply Lemma~\ref{lem:buffered-inverses} to $F_0=G_0$ and $F=G$.
The common mixed rectangles and gradient rectangles are chosen by that
lemma, with centers given by the model maps. They are the rectangles to
be used in the first B application: Lemma~\ref{lem:sharpB} uses their
source charts, and \eqref{eq:domain-provenance} places every resulting
normalized dual domain inside its $Q^-$. Further subdivisions, if any,
use restrictions of these branches; uniqueness on the original source
ball ensures agreement. Formula \eqref{eq:inverse-stability} with
$k=13$ proves the $C^{13}$ comparison from the available $C^{17}$ bound.
The gradient inverse $I=(\grad G)^{-1}$ has a uniform $C^{23}$ bound,
because differentiating $\grad G(I(y))=y$ to order $r\le23$ uses
derivatives of $G$ of order at most $r+1\le24$ and the uniformly bounded
inverse Hessian. Hence $\Phi=G\circ I-y\cdot I$ has, in particular,
a uniform $C^{22}$ bound on a neighborhood of $\overline{Q^+}$.

Here is the buffer and branch argument for the next A operation. Reserve
an intermediate rectangle $Q^\circ$ whose distance from $Q^-$ and from
the complement of $Q^+$ is at least a fixed $b_1>0$. Choose
$c_0<b_1/4$, and decrease it later if necessary. Then
\begin{equation}\label{eq:translation-buffer}
 \overline{Q^\circ}+[0,c_0]e_1\Subset Q^+.
\end{equation}
The ordered partition \eqref{eq:chart-partition} and the triangle
inequality are applied before A. For a single branch contribution with
normalized domain $E\subset Q^-$, a positive normalized shift
$\lambda e_1$ has correlation domain
\begin{equation}\label{eq:normalized-correlation-domain}
 E_\lambda=E\cap(E-\lambda e_1).
\end{equation}
Both endpoint phase values are evaluations of this same $\Phi$.
For $y\in E_\lambda$, the segment $y+[0,\lambda]e_1$ lies in the
convex rectangle $Q^-$, even if it leaves $E$. Formula
\eqref{eq:translation-buffer} additionally defines the same difference
quotient on a full neighborhood of $\overline{Q^\circ}$, without any
assumption that $E$ is connected or has interior. Thus no continuation
through a domain boundary or onto a different inverse branch is used.

For every spatial multi-index $|\alpha|\le9$, subtract $K_0$ in the
integral formula to obtain
\begin{align*}
 \partial^\alpha(\Phi_\lambda-K_0)(y)
 &=\int_0^1\partial^\alpha(\Phi_p-\Phi_{0,p})
                           (y+t\lambda e_1)\dd t\\
 &\quad+\int_0^1\bigl[\partial^\alpha\Phi_{0,p}
               (y+t\lambda e_1)-\partial^\alpha\Phi_{0,p}(y)\bigr]\dd t.
\end{align*}
The first integral is $O(e_0+\beta)$ by the $C^{13}$ comparison; the
second is $O(\lambda)$ by the uniform next derivative of $\Phi_0$.
This proves \eqref{eq:Phi-lambda-close}, with no
$(e_0+\beta)/\lambda$ loss. For $|\alpha|\le21$ the same formula,
without subtraction, and the $C^{22}$ bound of $\Phi$ give the uniform
$C^{21}$ bound for $\Phi_\lambda$ on $Q^\circ$.

Every point of $Q^\circ$ belongs to the chosen model inverse environment
of $G_0$, with its inverse in a fixed positive compact source set.
Lemma~\ref{lem:model} therefore supplies uniform determinant and pure
second margins for $K_0$ on a neighborhood of $\overline{Q^\circ}$.
They transfer to $\Phi_\lambda$ by its $C^9$ comparison, after decreasing
the three small upper bounds. Apply Lemma~\ref{lem:buffered-inverses}
again, now on $\overline{Q^-}$ with ambient domain $Q^\circ$, to
$F_0=K_0$, $F=\Phi_\lambda$. The number of source centers and their
buffers are uniform: apply the geometric construction after translating
the first-stage rectangles to fixed rectangles, and translate the source
charts back. Their side lengths and buffers are fixed, so a fixed finite
grid of offsets from the model centers suffices. Every resulting inverse
identity and phase is then written in the original $(p,q)$ coordinates.
These second-stage source charts are used in the second B
application, again with an ordered partition before estimating any sums.

On the common terminal rectangles, \eqref{eq:inverse-stability} with
$k=5$ and the $C^9$ comparison prove the stated $C^5$ closeness of
$\Psi$ and $\Psi_0$. The model fourth/fifth lower bound holds on these
whole compact inverse environments by Lemma~\ref{lem:model}. Choose
$e_0,\beta,c_0$ so that the two derivative errors are each less than
half its fixed margin; the same-coordinate maximum for $\Psi$ follows.
The derivative budget $C^{17}\to C^{13}\to C^9\to C^5$ thus has spare
derivatives at both inverse steps. Finally the $C^{21}$ bound on
$\Phi_\lambda$ and its determinant margin give a $C^{20}$ bound on
its gradient inverse. Its Legendre formula gives the required $C^{20}$
bound on $\Psi$. This proves \eqref{eq:uniform-high-regularity}.

For either $\mathcal H=G$ or $\mathcal H=\Phi_\lambda$, restoring the
physical scales consequently gives
\begin{equation}\label{eq:physical-derivative-bounds}
 \left|\partial_s^i\partial_t^j
       \bigl[P\mathcal H(s/U,t/V)\bigr]\right|
 \le C P U^{-i}V^{-j},\qquad i+j\le20.
\end{equation}
Thus both B applications in Proposition~\ref{prop:L5} have the required
upper bounds with constants independent of scale and aspect ratio.
The uniform semialgebraic assertions for these exact branches and all
subsequent sharp restrictions are verified separately in the next lemma.
\end{proof}

\begin{lemma}[Complexity transfer for the exact phase construction]
\label{lem:exact-complexity}
Fix the initial sharp-domain format and the finitely many derivative
orders and operations used in this proof. The exact families
$\mathcal F_{e_0},G,\Phi,\Phi_\lambda,\Psi$, their spatial derivatives
up to every specified order, all mixed and gradient inverse branches in
Lemma~\ref{lem:transfer}, and all transformed domains in
Lemma~\ref{lem:sharpB} have uniformly bounded semialgebraic format.
The assertion is joint in the parameters, including $e_0,\beta,\omega$,
$\lambda$, positive physical scales, real rectangle endpoints and
derivative thresholds. It remains uniform as $\beta$ or $\lambda$ tends
to zero through positive values. Their derivative extensions at zero
can be included as separate pieces.

In particular the A correlations, the original stationary restrictions,
and every individual terminal derivative layer have uniformly bounded
numbers of intervals and points in each coordinate fiber. No parameter
lower bound depending on $x$ is required for this complexity assertion.
\end{lemma}
\begin{proof}
An explicit graph description starts the construction. For $u,v>0$,
the equality $w=\mathcal F_{e_0}(u,v)$ is equivalent to the existence
of $r,c,t>0$ satisfying
\begin{equation}\label{eq:radical-graph}
 r^2=u,\qquad vc^3=v+e_0u,\qquad
 t^2(c^2+c+1)^3=27,\qquad vw=urt.
\end{equation}
Positivity selects every radical branch uniquely. These are fixed-degree
polynomial equations. The shifts add the polynomial relations
$u'=u+\beta\omega_1$, $v'=v+\beta\omega_2$ and
$\omega_1^2+\omega_2^2=1$. If $w_0,w_1$ are the two values from
\eqref{eq:radical-graph}, the graph of $G$ for $\beta>0$ adds
$\beta w=w_1-w_0$. At $\beta=0$ it is the directional derivative
$\omega\cdot\grad\mathcal F_{e_0}$, whose graph is supplied by
Lemma~\ref{lem:sa-calculus}. This proves uniform format without any
coefficient bound involving $1/\beta$.

Spatial derivative graphs of $G$ are obtained from that same lemma.
For a specified source ball $B_j$, the graph of its gradient inverse
is given by
\[
 z\in B_j,\qquad (p,q)=\grad G(z).
\]
Lemma~\ref{lem:buffered-inverses} proves uniqueness on the whole
buffered target rectangle. Adjoining the graph value $w=G(z)$ and the
polynomial relation $v=w-pz_1-qz_2$ gives the graph of $\Phi$.
The mixed inverse is handled by $p=G_1(z)$ and $q=z_2$, or by the
specified second pivot. All target rectangles have model graph values
as centers and fixed radii, and the pivot selection uses a fixed
derivative inequality. They therefore have bounded format as well.

For $\lambda>0$, the graph of $\Phi_\lambda$ uses two values $v_0,v_1$
of this same selected $\Phi$ and the equation $\lambda w=v_1-v_0$.
At zero use $w=\Phi_p$. Repeating the derivative-graph and specified
inverse-graph constructions gives the terminal inverse and $\Psi$.
The integrals in Lemma~\ref{lem:transfer} prove analytic norm estimates;
the graph argument here uses only the finite difference formulas.
Thus it does not require an assertion that integration preserves
semialgebraicity.

For domains, an A correlation is the single formula
$z\in D$, $z+r\in D$. Physical rescaling uses polynomial graph
relations such as $s=Uz_1$, $\xi U=Pp$, with positive scales; it does
not increase degrees as the scales vary. The two B transformations
use exactly the finitely many operations in
\eqref{eq:sharp-domain-family}, with every rectangle endpoint retained
as a free parameter. The initial stationary expressions
\eqref{eq:mnu}--\eqref{eq:wnu} are also compositions of rational
operations and positive radicals of fixed degrees, on their stated
positive frequency range. Their domain restrictions therefore fit the
same calculus. None of these steps inserts arithmetic support into a
real domain.

Finally, for example, a derivative layer is described by adjoining a
graph variable $w=\Psi_{2222}(y)$ and the inequalities
$\delta/2\le|w|\le\delta$, and projecting out $w$. Any fifth-derivative
or sign restriction uses the same rule. The number of dyadic levels
may grow, but each level is one member of this fixed family with free
parameter $\delta$; no growing union is used to obtain a component
bound. There is a fixed number of graph, projection and Boolean
operations along each construction path, and every derivative order is
fixed. Lemma~\ref{lem:sa-calculus} supplies a single resulting format
bound, followed by a single fiber-component bound. It includes coincident
endpoints, vanishing polynomial coefficients and the axis directions
of $\omega$. This proves all the assertions.
\end{proof}

\begin{remark}\label{rem:labels}
The frequencies retain their original labels: $p=G_u$, $q=G_v$.
Exchanging the order of the two B transforms changes only the order of
calculation. The next A shift always acts on $p$, so its limiting phase is
$\Phi_{0,p}=-u$. Replacing it by $\Phi_{0,q}$ would require a different
derivative calculation. The actual oscillating phase is never replaced by
its model in the sum; derivative closeness is the only approximation used.
\end{remark}

\section{Two seven-term local estimates}
\label{sec:local}

\begin{lemma}[Fourth-derivative sublevel decomposition]\label{lem:sublevel-fourth}
Let $\phi$ belong to a uniformly $C^5$-bounded semialgebraic family of
bounded description complexity on a fixed interval of length $O(1)$, and
suppose that
\[
 \max\bigl(|\phi^{(4)}(t)|,|\phi^{(5)}(t)|\bigr)\ge c>0.
\]
For $L\ge1$, $P>0$ and every interval restriction $I$ within that interval,
\begin{equation}\label{eq:sublevel-fourth}
 \left|\sum_{n:\,n/L\in I}e\bigl(P\phi(n/L)\bigr)\right|
 \ll 1+L^{5/7}P^{1/14}+L^{29/28}P^{-1/14}.
\end{equation}
The estimate also holds on a union of a bounded number of intervals and
points. Constants are uniform in the family and in all real endpoints.
\end{lemma}
\begin{proof}
The case $L=O(1)$ is trivial. On $|\phi^{(4)}|\ge c/2$, split into its
bounded number of interval components and use \eqref{eq:fourth-test}
with ambient length $L$ and $\lambda\asymp P/L^4$. Isolated points in
this part contribute $O(1)$.

On the remaining set, $|\phi^{(5)}|\ge c$. Decompose it into dyadic
layers
\[
 \delta<|\phi^{(4)}|\le2\delta,\qquad
 L^{-1}\lesssim\delta\lesssim1,
\]
and a final set $|\phi^{(4)}|\lesssim L^{-1}$. Each layer has a uniformly
bounded number of interval and point components. On each interval the
fifth derivative has constant sign. The mean value theorem therefore
bounds its normalized length by $O(\delta)$. In the physical variable,
the available interval length and the fourth derivative size satisfy
\[
 M\ll\delta L,\qquad \lambda\asymp P\delta/L^4.
\]
Here $\delta L\gtrsim1$, so a suitable fixed multiple of $\delta L$
also bounds the number of integer points. Applying \eqref{eq:fourth-test}
gives
\begin{equation}\label{eq:sublevel-layer}
 \delta^{15/14}L^{5/7}P^{1/14}
 +\delta^{19/28}L^{29/28}P^{-1/14}.
\end{equation}
Both powers of $\delta$ are positive, and the dyadic sums converge.
Isolated points in such a layer can be included in the same bound:
the sum of the two terms in the fourth-derivative estimate is at least
$2M^{7/8}$ for $M\ge1$, by the arithmetic--geometric mean inequality.

Each component of the final set has normalized length $O(L^{-1})$,
again by $|\phi^{(5)}|\ge c$ and the mean value theorem. It contributes
$O(1)$ points, and there are $O(1)$ components. Semialgebraic finiteness
provides all the component bounds uniformly, also after intersecting
with $I$. This proves \eqref{eq:sublevel-fourth}.
\end{proof}

\begin{proposition}\label{prop:L5}
Let $1\le H\le X$, let $e_0=H/X$ be sufficiently small, and let $\beta>0$
be sufficiently small. Let $D$ be a permitted sharp domain in an $H$ by $X$
box whose normalized coordinates lie in the fixed positive compact set.
For $Z>0$ put
\[
 S_5=\sum_{(h,n)\in D\cap\Z^2}
       e\bigl(ZG_{e_0,\beta,\omega}(h/H,n/X)\bigr).
\]
Uniformly over unit directions $\omega$, domains and sharp endpoints,
both of the following estimates hold:
\begin{align}
 |S_5|\ll x^\eta\bigl(&Z^{13/25}H^{1/2}X^{1/5}
       +Z^{89/200}H^{1/2}X^{13/40}+\sqrt{XZ}\notag\\
       &+Z^{6/25}H^{1/2}X^{2/5}
       +ZH^{-1/4}+X+HX/\sqrt Z\bigr),\label{eq:L5}\\
 |S_5|\ll x^\eta\bigl(&Z^{13/25}H^{21/50}X^{7/25}
       +Z^{89/200}H^{101/175}X^{347/1400}+\sqrt{XZ}\notag\\
       &+Z^{6/25}H^{27/50}X^{9/25}
       +ZH^{-1/4}+X+HX/\sqrt Z\bigr).\label{eq:L5safe}
\end{align}
We call these the long-direction and conservative bounds, respectively.
The scales have at most polynomial size in $x$; equivalently one may
retain the fixed logarithmic losses in the proof.

For later reference, their powers $(a_j^\star,b_j^\star,d_j^\star)$ in
$Z^{a_j^\star}H^{b_j^\star}X^{d_j^\star}$ are
\begin{equation}\label{eq:local-powers}
\begin{array}{c|rrr|rrr}
 &\multicolumn{3}{c|}{\star=\mathrm L}
 &\multicolumn{3}{c}{\star=\mathrm S}\\
 j&a_j^{\mathrm L}&b_j^{\mathrm L}&d_j^{\mathrm L}
  &a_j^{\mathrm S}&b_j^{\mathrm S}&d_j^{\mathrm S}\\ \hline
 1&13/25&1/2&1/5&13/25&21/50&7/25\\
 2&89/200&1/2&13/40&89/200&101/175&347/1400\\
 3&1/2&0&1/2&1/2&0&1/2\\
 4&6/25&1/2&2/5&6/25&27/50&9/25\\
 5&1&-1/4&0&1&-1/4&0\\
 6&0&0&1&0&0&1\\
 7&-1/2&1&1&-1/2&1&1
\end{array}.
\end{equation}
\end{proposition}
\begin{proof}
If $Z<1$, the last term bounds the trivial $O(HX)$ count. If $1\le Z<X$,
Lemma~\ref{lem:sharpB} and \eqref{eq:trivial-dual}, with lengths $H,X$,
give the trivial dual main term $Z+H+X+HX/Z$. Together with
$\mathcal R(H,X,Z)$, this is bounded by
\[
 (X+\sqrt{XZ}+HX/\sqrt Z)\log(2+H+X+Z).
\]
This treats every short initial dual range.

Suppose $Z\ge X$, and set
\begin{equation}\label{eq:UVW}
 U=Z/H,\qquad V=Z/X,\qquad W=HX/Z,\qquad U\ge V\ge1.
\end{equation}
The first two-dimensional B inequality gives $W$ times a maximal dual sum
$S_6$, plus $\mathcal R(H,X,Z)$. Fix one chart and one actual dual domain
$E$ before applying A. Its exact phase is $Z\Phi(k/U,l/V)$, with the
frequency labels of Remark~\ref{rem:labels}. A positive integer shift
$\tau$ in the first coordinate has exact correlation phase
\[
 Z\bigl(\Phi((k+\tau)/U,l/V)-\Phi(k/U,l/V)\bigr)
 =P\Phi_\lambda(k/U,l/V),\quad
 P=H\tau,\quad\lambda=\tau/U.
\]
Lemma~\ref{lem:transfer} applies when $\lambda\le c_0$.

Let $T(\tau)$ denote this unweighted correlation sum. Apply the sharp B
inequality again, with lengths $U,V$ and parameter $P$. If $P\ge U$,
both terminal dual lengths
\[
 R_1=P/U,\qquad R_2=P/V,\qquad R_1\le R_2
\]
are at least one. The exact terminal phase is
$P\Psi(\xi/R_1,\zeta/R_2)$, where $\Psi=\Leg\Phi_\lambda$.
For each fixed first coordinate, the $C^5$ bounds and the
same-coordinate fourth/fifth derivative condition in
Lemma~\ref{lem:transfer} permit Lemma~\ref{lem:sublevel-fourth} along
the second coordinate. Lemma~\ref{lem:exact-complexity} gives a uniform component bound for every
actual terminal fiber and every derivative layer. Consequently its sum is bounded by
\[
 1+R_2^{5/7}P^{1/14}+R_2^{29/28}P^{-1/14}.
\]
Counting the $O(R_1)$ first-coordinate frequencies and multiplying by
$UV/P$ gives
\[
 V+P^{15/14}R_2^{-2/7}+P^{13/14}R_2^{1/28}
 =V+P^{11/14}V^{2/7}+P^{27/28}V^{-1/28}.
\]
The $V$ term is already included in $\mathcal R(U,V,P)$.
All derivative-level restrictions occur after the last B transform.
They only truncate the actual terminal fibers. Smooth amplitudes have
already been removed by rectangular Abel summation, so no derivatives
of the layer endpoints or additional stationary boundary errors arise.

If $P<U$, use the complete trivial dual expression
$P+U+V+UV/P$. Since $P\le U$, it is covered by $\mathcal R(U,V,P)$.
Thus in every terminal range,
\begin{equation}\label{eq:terminal}
 |T(\tau)|\ll x^\eta\left(
 P^{11/14}V^{2/7}+P^{27/28}V^{-1/28}
 +\mathcal R(U,V,P)\right).
\end{equation}
As $V\le U$, this also holds if the first $V^{2/7}$ is enlarged to
$U^{2/7}$. This gives the conservative alternative without imposing
any additional derivative hypothesis.

Choose $\star\in\{\mathrm L,\mathrm S\}$, put
$B_{\mathrm L}=V$, $B_{\mathrm S}=U$, and apply the one-axis version of
Lemma~\ref{lem:A} with
\begin{align}
 Q_{0,\mathrm L}&=Z^{24/25}H^{-1}X^{-2/5},&
 Q_{0,\mathrm S}&=Z^{24/25}H^{-21/25}X^{-14/25},\notag\\
 Q&=\max\{1,\lfloor\min(Q_{0,\star},c_0U)\rfloor\}.
 \label{eq:Q-choice}
\end{align}
Extend $T(\tau)$ to negative shifts by its sharp-domain correlation
definition. Translating the summation variable gives
$T(-\tau)=\overline{T(\tau)}$, so both signs cost only a factor of two.
For $Q\ge2$, every $1\le\tau<Q$ satisfies $\tau/U<c_0$.
Average positive powers of $\tau$ using $Q$ and bound negative powers
conservatively by $\tau\ge1$. Then \eqref{eq:terminal} yields
\begin{align}
 |S_6|^2\ll x^\eta\biggl\{\frac{(UV)^2}{Q}
 +UV\biggl(& (HQ)^{11/14}B_\star^{2/7}
       +(HQ)^{27/28}V^{-1/28}+U+V\notag\\
       &+\sqrt{HQ}+\frac{UV}{\sqrt H}+\frac{UV}{H}
       \biggr)\biggr\}.\label{eq:S6-square}
\end{align}
When $Q=1$, there are no correlations and $|S_6|\ll UV$ suffices.
In either branch,
\begin{equation}\label{eq:Q-diagonal}
 \frac Z{\sqrt Q}\ll\frac Z{\sqrt{Q_{0,\star}}}+\sqrt{HZ}
 \le\frac Z{\sqrt{Q_{0,\star}}}+\sqrt{XZ}.
\end{equation}
For $Q\ge2$, use comparability of the floor to the minimum. For $Q=1$,
either $Q_{0,\star}<2$ or $c_0U<2$, and one of the right-hand terms
bounds $Z$.

Since $W\sqrt{UV}=\sqrt{HX}$, the diagonal contributes $Z/\sqrt Q$.
The first off-diagonal term contributes, respectively,
\begin{align*}
 A_{\mathrm L}&=Z^{1/7}H^{25/28}X^{5/14}Q^{11/28},\\
 A_{\mathrm S}&=Z^{1/7}H^{3/4}X^{1/2}Q^{11/28}.
\end{align*}
The remaining off-diagonal contributions are common to both choices:
\[
 Z^{-1/56}H^{55/56}X^{29/56}Q^{27/56},\quad
 \sqrt{XZ},\quad\sqrt{HZ},\quad
 H^{3/4}X^{1/2}Q^{1/4},\quad
 ZH^{-1/4},\quad ZH^{-1/2}.
\]
In the $Q\ge2$ branch only, replace nonnegative powers of $Q$ by
the corresponding powers of $Q_{0,\star}$. The choice
$Q_{0,\star}$ balances the diagonal with $A_\star$ and gives the first
term of the corresponding bound. The next contributions give the
second, third and fourth terms. Since $H\ge1$, $ZH^{-1/2}$ is covered
by $ZH^{-1/4}$. Finally, for $Z\ge X\ge H\ge1$,
\[
 \mathcal R(H,X,Z)\ll
 (HX/\sqrt Z+X+\sqrt{XZ})\log(2+H+X+Z).
\]
For $Q=1$, use only \eqref{eq:Q-diagonal}; no off-diagonal bound is
needed. This proves both \eqref{eq:L5} and \eqref{eq:L5safe} for every
$Z>0$ and every allowed sharp restriction.
\end{proof}

\section{The original divisor sum and the Case A reduction}
\label{sec:reduction}
Choose $Y=\lfloor x^\rho\rfloor$ in \eqref{eq:voronoi}. Partition each of
the three factors of $n$ into dyadic intervals $[2^j,2^{j+1})$, $j\ge0$.
In a fixed factor box, designate the largest scale by $X$, breaking ties
by a fixed index order, and call the other scales $A,B$. Put
\[
 M=AB,\qquad N=ABX=MX.
\]
Every nonempty cutoff box has
\begin{equation}\label{eq:box-range}
 N\le Y,\qquad N^{1/3}\le X\le N,
 \qquad N\le abn<8N.
\end{equation}
Define
\[
 a_m=\#\{(a,b)\in I_A\times I_B:ab=m\},
 \qquad a_m\le d(m),\qquad \sum_{m\in I_M}|a_m|^2\ll Mx^\eta,
\]
where $I_M$ is the full ambient product interval, of length $O(M)$.
No pointwise ordering of the factors is imposed.

One-dimensional Abel summation in the product removes its weight
$n^{-2/3}$ at cost $O(N^{-2/3})$: the normalized weight has bounded
variation on $[N,8N)$. Thus it suffices to prove, uniformly in real product
cutoffs $Y_1,Y_2$, that
\begin{align}
 S_1&=\sum_{m\in I_M}a_m
       \sum_{\substack{n\in I_X\\Y_1<mn\le Y_2}}
             e\bigl((Rmn)^{1/3}\bigr),\label{eq:S1}\\
 x^{1/3}N^{-2/3}|S_1|&\ll x^{\theta+\eta}.
       \label{eq:box-target}
\end{align}
The physical scale is \eqref{eq:F-convention}. The number of factor boxes
is $O(\log^3x)$. Boxes with $N\le x^{u_0}$ are trivial; their contribution
is included explicitly in Section~\ref{sec:assembly}.

\subsection{Cauchy and the signed first A}
For the moment let $1\le q\ll X$ be an integer. Set
\[
 z_m(n)=1_D(m,n)e\bigl((Rmn)^{1/3}\bigr),\qquad
 D=\{(m,n)\in I_M\times I_X:Y_1<mn\le Y_2\},
\]
with zero extension in $n$. Cauchy gives
\begin{equation}\label{eq:full-Cauchy}
 |S_1|^2\le\left(\sum_{m\in I_M}|a_m|^2\right)
        \sum_{m\in I_M}\left|\sum_n z_m(n)\right|^2.
\end{equation}
The second sum runs over the full $m$ interval, including $a_m=0$.
Restricting it to the arithmetic support of $a_m$ would not justify the
next B transform.

There is a common interval of $K_q\ll X+q$ integer points containing the
support of $\sum_{r=1}^q z_m(n+r)$ for every $m$. Define
\[
 C_m(h)=\sum_n z_m(n+h)\overline{z_m(n)},\qquad
 C(h)=\sum_{m\in I_M}C_m(h).
\]
The exact shift identity and expansion of its squared norm give
\begin{equation}\label{eq:signed-A}
 |S_1|^2\le
 \left(\sum_{m\in I_M}|a_m|^2\right)\frac{K_q}{q}
 \left[C(0)+2\Re\sum_{1\le h<q}(1-h/q)C(h)\right].
\end{equation}
The bracket is nonnegative, being a sum of squared norms divided by $q$.
Moreover $C(0)\ll N$. On each dyadic interval $I_H=[H,2H)$, the endpoint
size plus variation of $1-h/q$ on $I_H\cap[1,q)$ is at most two. Abel
summation in $h$ therefore gives
\begin{equation}\label{eq:A0}
 |S_1|^2\ll x^\eta\left(\frac{N^2}{q}
             +\frac Nq\sum_{\substack{H\ \mathrm{dyadic}\\H\le q}}S_2(H)\right),
\end{equation}
where
\[
 S_2(H)=\sup_{J\subset I_H\cap[1,q)}
            \left|\sum_{h\in J}C(h)\right|,
\]
and $J$ ranges over real intervals. It is essential that the joint signed
$h$ sum is kept until this point; there is no extra factor $H$ from
removing the triangular weight. If $q=1$, the correlation sum is empty.

\subsection{The first stationary phase, with its exact weight}
Fix $h\asymp H>0$ and $n\asymp X$, and set
\begin{equation}\label{eq:T-L}
 d_h(n)=(n+h)^{1/3}-n^{1/3}>0,
 \quad T=FH/X,\quad L=T/M=FH/N,\quad \mathcal A=M/\sqrt T.
\end{equation}
For each such pair, the $m$ phase and domain are
\[
 g(m)=R^{1/3}d_h(n)m^{1/3},
\]
\[
 I_M\cap(Y_1/n,Y_2/n]\cap(Y_1/(n+h),Y_2/(n+h)],
 \qquad n,n+h\in I_X.
\]
The domain is an interval without an arithmetic restriction. The curvature
has size $T/M^2$ and the next two derivatives obey the B hypotheses.

The stationary point, phase and weight are exactly
\begin{align}
 m_\nu&=\frac{\sqrt R\,d_h(n)^{3/2}}{(3\nu)^{3/2}},
       \label{eq:mnu}\\
 f_\nu(h,n)&=g(m_\nu)-\nu m_\nu
     =\frac{2\sqrt R}{3\sqrt{3\nu}}d_h(n)^{3/2},
       \label{eq:fnu}\\
 g''(m_\nu)&=-\frac{2\nu}{3m_\nu},\qquad
 w_\nu=\frac{R^{1/4}d_h(n)^{3/4}}
                  {\sqrt2\,3^{1/4}\nu^{5/4}}.
       \label{eq:wnu}
\end{align}
The Maslov factor is the constant $e(-1/8)$. Every stationary frequency
is a positive integer with $cL\le\nu\le CL$, for fixed $c,C>0$.
Either no integer frequency occurs, or $L\ge1/C$ and the number of
frequencies is $O(L)$. This is why a separate $+1$ in the frequency count
causes no loss; the zero frequency is absent.

Original endpoint changes and isolated fibers cost $O(1)$ per $(h,n)$;
dual endpoint conventions cost $O(\mathcal A)$ per pair. Charge them
before grouping by $\nu$. Together with the B error this gives
\begin{equation}\label{eq:Err0}
 E_0\ll \bigl(HX+N\sqrt{HX/F}\bigr)\log x.
\end{equation}
Thus an endpoint half-weight never becomes a discontinuous coefficient
to which smooth Abel summation is applied.

For explicit normalization, write $h=Hu$, $n=Xv$, $\nu=L\ell$, and put
\[
 e_0=H/X,\qquad c_e=(1+e_0u/v)^{1/3},\qquad
 \kappa_e=\frac3{c_e^2+c_e+1}.
\]
Then
\begin{align}
 d_h(n)&=\frac{H}{3X^{2/3}}u v^{-2/3}\kappa_e,\notag\\
 f_\nu(h,n)&=\frac2{27}T\ell^{-1/2}\mathcal F_{e_0}(u,v)
            =T_\nu\mathcal F_{e_0}(u,v),\label{eq:normalized-firstB}\\
 T_\nu&=\frac{2\sqrt R H^{3/2}}{27\sqrt\nu X}\asymp FH/X,\notag\\
 w_\nu&=\frac{\mathcal A}{3\sqrt2}
       \ell^{-5/4}u^{3/4}v^{-1/2}\kappa_e^{3/4}.\notag
\end{align}
Since $\ell$ stays in a fixed positive compact interval, $w_\nu/\mathcal A$
has uniformly bounded normalized $C^2$ norm on a larger positive rectangle,
before imposing any stationary domain or product cutoff.

For a fixed endpoint convention the stationary domain may be written as
\begin{align*}
 D_{\nu,J,Y}=\{(h,n):
 &h\in J,\ n,n+h\in I_X,\ m_\nu\in I_M,\\
 &Y_1<n m_\nu,\quad (n+h)m_\nu\le Y_2\}.
\end{align*}
It has bounded semialgebraic complexity. Lemma~\ref{lem:Abel} bounds its
weighted sum by $\mathcal A\mathcal M_3(\nu)$, where
\[
 \mathcal M_3(\nu)=
 \sup_{J,K_h,K_n}
 \left|\sum_{D_{\nu,J,Y}\cap(K_h\times K_n)}e(f_\nu(h,n))\right|.
\]
All cuts are real intervals. The order of the maxima and frequency sums is
\begin{align*}
 S_2(H)&\le
 \sup_J\sum_\nu\left|\sum_{D_{\nu,J,Y}}w_\nu e(f_\nu)\right|+E_0\\
 &\le\sum_\nu\sup_J
         \left|\sum_{D_{\nu,J,Y}}w_\nu e(f_\nu)\right|+E_0\\
 &\ll\mathcal A\sum_{\nu\asymp L}\mathcal M_3(\nu)+E_0.
\end{align*}
Since $\mathcal A L=\sqrt T$, we obtain
\begin{equation}\label{eq:A1}
 S_2(H)\ll x^\eta\left(
       \sqrt{FH/X}\sup_\nu\mathcal M_3(\nu)
       +N\sqrt{HX/F}+HX\right).
\end{equation}
An empty frequency supremum contributes zero. Every $\nu$ may have its own
maximizing rectangles. The bound is uniform in their endpoints, and no
regularity of these endpoints as functions of $\nu$ is required.

\subsection{The second A and the complete shift average}
Fix one unweighted sum $S_3$ occurring in $\mathcal M_3(\nu)$. Apply
Lemma~\ref{lem:A} with lengths $H,X$. In a full two-axis choice suppose
\[
 q_1\asymp\sqrt{\Pi H/X},\quad
 q_2\asymp\sqrt{\Pi X/H},\quad
 \pint=q_1q_2\asymp\Pi,\quad
 \alpha=\sqrt{\Pi/(HX)}.
\]
In an axis choice take $q_1=1$, $q_2\asymp\Pi$ and $\alpha=\Pi/X$.
The exact choices and all admissibility conditions are verified in
Section~\ref{sec:caseA}.

For each nonzero shift $r=(r_h,r_n)$, set
\[
 \beta=\sqrt{(r_h/H)^2+(r_n/X)^2},\qquad
 \omega=\beta^{-1}(r_h/H,r_n/X),\qquad Z=T_\nu\beta.
\]
Its correlation is an unweighted sum over the exact
$D_3\cap(D_3-r)$ with phase $ZG_{e_0,\beta,\omega}$. Consequently
Proposition~\ref{prop:L5} applies once the smallness of the normalized
shifts is established. No arbitrary coefficient extension of that
proposition is needed.

For every $a\ge0$, the average of $\beta^a$ is $O(\alpha^a)$. The
negative-power average must also be kept. Direct summation gives
\begin{align}
 \frac1\pint\sum_{r_n\ne0}\beta^{-1/2}
       &\ll\sqrt{X/q_2},\label{eq:negative-average}\\
 \frac1\pint\sum_{\substack{r_n=0\\r_h\ne0}}\beta^{-1/2}
       &\ll\frac{\sqrt{H/q_1}}{q_2}.\notag
\end{align}
Indeed use $\beta\ge |r_n|/X$ in the first sum and
$\beta\ge |r_h|/H$ in the second, followed by
$\sum_{j<q}j^{-1/2}\ll\sqrt q$. In the full case both contributions are
$O(\alpha^{-1/2})$. If $q_1=1$, the second sum is empty; in the axis
case the first alone gives the same result. Nonnegative powers follow
from $\beta\ll\alpha$ and the $O(\pint)$ number of shifts.

Choose either $\star=\mathrm L$ or $\star=\mathrm S$ uniformly for
the fixed global parameters, before averaging the nonzero second shifts.
Averaging the seven terms of \eqref{eq:L5} or \eqref{eq:L5safe} yields
\begin{equation}\label{eq:A2}
 |S_3|^2\ll x^\eta\left(\frac{(HX)^2}{\Pi}
                       +HX\sum_{j=1}^7\mathcal J_j\right),
\end{equation}
where, for the selected row of \eqref{eq:local-powers},
\begin{equation}\label{eq:Jterms}
 \mathcal J_j=
 \left(\frac{FH\alpha}{X}\right)^{a_j^\star}
 H^{b_j^\star}X^{d_j^\star},\qquad 1\le j\le7.
\end{equation}
The comparability $T_\nu\asymp FH/X$ is uniform, including for the
negative power. In both alternatives,
\[
 \mathcal J_7=\frac{HX}{\sqrt{FH\alpha/X}}.
\]
It is the average of a retained boundary error. Replacing $\beta$ by its
smallest possible value would not give the stated estimate.
The choice between the two bounds may depend on $N,H,X$, but not on
an unaccounted exceptional set of individual shifts.

\section{Case A: parameter choices and their certification}
\label{sec:caseA}
Assume
\begin{equation}\label{eq:caseA-range}
 x^{u_0}\le N\le x^\rho,\qquad
 N^{1/3}\le X\le N^{9/20}.
\end{equation}
In this section and Appendix~\ref{app:affine}, the letters $u,v,h,f$
denote logarithmic exponents:
\[
 N=x^u,\quad X=x^v,\quad H=x^h,\quad
 \Fzero=x^f,\quad f=(1+u)/3.
\]
They are not the normalized positive coordinates used for the local models.

Choose
\begin{equation}\label{eq:first-q}
 q_*=(N/x^{u_0})^{2/3},\qquad
 q=\max\{1,\lfloor q_*\rfloor\},\qquad
 \hcap=\frac23(u-u_0).
\end{equation}
Throughout \eqref{eq:caseA-range},
\begin{equation}\label{eq:initial-cap-margin}
 v-\hcap=\frac{11}{1209}+\frac{\rho-u}{3}
                  +\left(v-\frac u3\right).
\end{equation}
Thus $q_*/X\le x^{-11/1209}$ and the first A is admissible. If $q=1$,
$q_*<2$ implies $N=O(x^{u_0})$, so the trivial bound already gives
\eqref{eq:box-target}. We henceforth treat the nonempty correlations.

\subsection{The complete small-\texorpdfstring{$H$}{H} range}
Put
\begin{equation}\label{eq:Ht}
 H_{\mathrm t}=N^{2/3}\Fzero^{-1/3}X^{-1/3},\qquad
 \hthreshold=(2u-f-v)/3.
\end{equation}
The following exact identities hold on the whole Case A range:
\begin{align}
 \hthreshold&=\frac{1007}{36270}+\frac{73}{180}(u-u_0)
                   +\frac13\left(\frac{9u}{20}-v\right),\notag\\
 f-\hthreshold-v&=\frac{733}{7254}+\frac{47}{90}(\rho-u)
                   +\frac23\left(\frac{9u}{20}-v\right),
                   \label{eq:Ht-margins}\\
 u-\hthreshold-v&=\frac{224}{1395}+\frac{13}{90}(u-u_0)
                   +\frac23\left(\frac{9u}{20}-v\right).\notag
\end{align}
In particular $H_{\mathrm t}$ grows as a fixed power of $x$, and
$H_{\mathrm t}X$ is smaller than both $\Fzero$ and $N$ by fixed powers.
For $H\le H_{\mathrm t}$, use the trivial bound
$\mathcal M_3(\nu)\ll HX$ in \eqref{eq:A1}. Its main term is
\[
 \sqrt{FH/X}\,HX=H^{3/2}\sqrt{FX}\ll N,
\]
and both errors in \eqref{eq:A1} are $O(N)$ by
\eqref{eq:Ht-margins}. This includes empty or short frequency sets. If
$H_{\mathrm t}>q_*$, all correlations are covered. Every remaining scale
lies in the exact region
\begin{equation}\label{eq:largeH-region}
 u_0\le u\le\rho,\qquad u/3\le v\le9u/20,
 \qquad \hthreshold\le h\le\hcap.
\end{equation}

\subsection{Nominal and integer second shifts}
Set
\begin{equation}\label{eq:c-p0-pb}
 c=75h+16v-26f,\qquad
 \mathfrak p_0=\frac{12h+79v-26f}{63},\qquad
 \mathfrak p_b=11h+2f+7v-8u.
\end{equation}
On $c\le0$ take the axis choice
\begin{equation}\label{eq:axis-choice}
 \Pi=H^{-1/76}X^{23/19}\Fzero^{-13/38},\qquad
 q_1=1,\quad q_2=\max\{1,\lfloor\Pi\rfloor\},\quad\alpha=\Pi/X.
\end{equation}
On $c\ge0$ take
\begin{align}
 \Pi_0&=H^{4/21}X^{79/63}\Fzero^{-26/63},\qquad
 \Pi_b=H^{11}\Fzero^2X^7N^{-8},\notag\\
 \Pi&=\max(\Pi_0,\Pi_b),\qquad
 r_1=\sqrt{\Pi H/X},\quad r_2=\sqrt{\Pi X/H},\label{eq:full-choice}\\
 q_j&=\max\{1,\lfloor r_j\rfloor\},\qquad
 \alpha=\sqrt{\Pi/(HX)}.\notag
\end{align}
Use the long-direction bound \eqref{eq:L5} on the axis region and on
the full region $\mathfrak p_0\ge\mathfrak p_b$. On the remaining full
region $\mathfrak p_b\ge\mathfrak p_0$, use the conservative bound
\eqref{eq:L5safe}. Thus the choice of $\star$ in \eqref{eq:Jterms}
is fixed before averaging the second shifts. At a common boundary either
of the applicable, certified choices may be used.

This distinction is needed for the secondary term. If the long-direction
seven terms were used throughout the boundary region with
$\mathfrak p=\mathfrak p_b$, its second deficit at
\[
 (u,v,h)=\left(\frac{265}{403},\frac{477}{1612},\frac{254}{1209}\right)
\]
would be $-4213/1934400$. The conservative choice makes every required
deficit nonnegative, as certified in Appendix~\ref{app:affine}.

The lower choice $\Pi_b$ is required by the last term in
\eqref{eq:Jterms}: in the full case its contribution to \eqref{eq:A1}
after \eqref{eq:A2} is comparable to
\begin{equation}\label{eq:Pi-boundary-term}
 H^{11/8}\Fzero^{1/4}X^{7/8}\Pi^{-1/8}.
\end{equation}
Thus $\Pi\ge\Pi_b$ pays this term at the exponent level. Using only
$\Pi_0$ would leave an uncovered part of the parameter region.

\subsection{What the exact certificates prove}
Appendix~\ref{app:affine} gives nonnegative rational combinations for every
required affine inequality. The three Case A regions are the axis region,
the full region with $\mathfrak p_0\ge\mathfrak p_b$, and the full region
with $\mathfrak p_b\ge\mathfrak p_0$. The 44 certificates on these regions
prove that every term of \eqref{eq:A2}, after taking a square root and
multiplying by $\sqrt{FH/X}$, is $O(N)$, and that both errors of
\eqref{eq:A1} are $O(N)$. They also prove
\begin{equation}\label{eq:A-admissibility}
 \Pi\ge x^{1/1000},\qquad
 \alpha\le x^{-1/1000},\qquad
 \Fzero H/N\ge x^{1/1000},\qquad H/X\le x^{-11/1209}.
\end{equation}
The stronger bound on $H/X$ also uses \eqref{eq:initial-cap-margin}
and $h\le\hcap$. These are identities throughout the regions, not
numerical tests at selected points.

For clarity, check that the actual integer shifts satisfy all local
hypotheses. In the full case,
\[
 \log_x\sqrt{\Pi_0H/X}=c/126\ge0,
\]
so $r_2\ge r_1\ge1$ and
\begin{equation}\label{eq:integer-product}
 r_j/2\le q_j\le r_j,\qquad \Pi/4\le\pint\le\Pi.
\end{equation}
The normalized lengths are at most $\alpha$, which tends to zero, and
$q_2\ge\sqrt\Pi/2$ grows. In the axis case $q_2\asymp\Pi$ for large $x$,
while the only first-coordinate correlation shift is zero. Thus the
negative-power averages \eqref{eq:negative-average} hold for the rounded
integers, and the nonnegative-power averages do also. The smallness of
$H/X$ and of every occurring $\beta$ follows from
\eqref{eq:A-admissibility}. The positive frequency scale condition follows
from its third inequality. The final local clamp in \eqref{eq:Q-choice}
uses the fixed $c_0$ already allowed in Proposition~\ref{prop:L5}.

Consequently \eqref{eq:A2} is legitimately applicable with these choices,
including its negative $Z$ power. Replacing $F$ by $3\Fzero$ changes only
fixed constants in every propagated term. The parameter cuts include their
boundaries and leave no gap. We have therefore proved
\[
 S_2(H)\ll Nx^\eta
\]
for every initial dyadic correlation scale, using either the small-$H$
argument or the certificates. Equation~\eqref{eq:A0} gives
$|S_1|\ll Nq^{-1/2}x^\eta$ after absorbing logarithms. Since $q\asymp q_*$
outside the already treated trivial branch, its weighted exponent is
\begin{equation}\label{eq:caseA-final}
 \frac13+\frac u3-\frac{\hcap}{2}
       =\frac13+\frac{u_0}{3}=\theta.
\end{equation}
This proves \eqref{eq:box-target} throughout Case A, uniformly in all
original product cutoffs.

\subsection{The balance determining the exponent}
The parameter choices can also be read from the active balanced box.
On $v=u/3$, balancing the first long-direction monomial with the second
A diagonal gives $\mathfrak p=\mathfrak p_0$. Its propagated exponent
in $S_2$ is
\[
 \frac{89f+177h-16v}{126}
 =\frac{89+73u+531h}{378}.
\]
Equating this to $u$ gives the first-A balance
$h_{\rm opt}=(305u-89)/531$. After the Voronoi weight, the corresponding
box exponent is
\[
 E_{\mathrm L}(u)=\frac13+\frac u3-\frac{h_{\rm opt}}2
                =\frac{443+49u}{1062}.
\]
The equality $E_{\mathrm L}(\rho)=2/3-\rho/3$ gives
$\rho=265/403$ and $\theta=541/1209$. This explains the constants;
the certificates, not this calculation alone, establish the entire range.

\section{Case B: longer largest factors}
\label{sec:caseB}
We now assume
\begin{equation}\label{eq:caseB-range}
 x^{u_0}\le N\le x^\rho,\qquad N^{9/20}\le X\le N,
 \qquad M=N/X,\quad L=F/X.
\end{equation}

\subsection{The first B transform and removal of its weight}
Apply the one-dimensional B transform to the $n$ sum in \eqref{eq:S1}
before any Cauchy inequality. Its stationary formulas are
\begin{align}
 n_l&=\frac{\sqrt{Rm}}{(3l)^{3/2}},\qquad
 f_0(m,l)=\frac{2\sqrt{Rm}}{3\sqrt{3l}},\notag\\
 w(m,l)&=\frac{(Rm)^{1/4}}{\sqrt2\,3^{1/4}l^{5/4}}.
              \label{eq:caseB-firstB}
\end{align}
Stationary frequencies are positive and $l\asymp L$. The endpoint and B
errors, summed against $|a_m|$, are
\begin{equation}\label{eq:caseB-first-error}
 \ll x^\eta(N/\sqrt F+M).
\end{equation}
If $X>F$, the stationary interval contains $O(1)$ integer frequencies or
none. Its main term is bounded by the same quantity, so
\begin{equation}\label{eq:caseB-short-first}
 |S_1|\ll x^\eta(N/\sqrt F+M)\qquad (X>F).
\end{equation}

Suppose $X\le F$, so $L\ge1$. At $m=Mu$, $l=Lv$, the exact weight is
\[
 w(m,l)=\frac{X/\sqrt F}{\sqrt2\,3^{1/4}}u^{1/4}v^{-5/4}.
\]
The normalized amplitude has bounded $C^2$ norm on a positive ambient
rectangle. Rectangular Abel removes it, retaining $a_m$ in the base array
and retaining the exact stationary domain, including the product cutoffs.
For each resulting unweighted-in-$l$ sum, Cauchy over the full $m$ interval
removes $a_m$, and A in $l$ gives unit-coefficient correlations on the
intersected and translated domain. The squared scale factors are
\[
 (X/\sqrt F)^2ML=N,\qquad
 (X/\sqrt F)^2(ML)^2=M^2F.
\]
Consequently
\begin{equation}\label{eq:caseB-A}
 |S_1|^2\ll x^\eta\left(
 \frac{M^2F}{q}+\frac Nq\sum_{1\le h<q}B_2(h)
          +\frac{N^2}{F}+M^2\right).
\end{equation}
Each $B_2(h)$ denotes the absolute value of an unweighted correlation,
uniformly over the resulting domains. For $q=1$ the correlations are
empty. Whenever $q\ge2$, impose
\begin{equation}\label{eq:caseB-cap}
 q\le\frac{\min(X,L)}{C\log x},
\end{equation}
where $C\ge1$ is fixed. In particular $q\le L$ and $h/L$ is small for
every nonempty correlation.

\subsection{The exact local correlation and both short-dual cases}
The correlation phase is
\begin{equation}\label{eq:caseB-correlation}
 f_1(m,l)=\frac{2\sqrt{Rm}}{3\sqrt3}
          \bigl((l+h)^{-1/2}-l^{-1/2}\bigr).
\end{equation}
Set $P=hX$ and $\delta=h/L$. Its exact normalized phase divided by $P$
is
\[
 \Gamma_\delta(u,v)=\frac{2\sqrt u}{3\sqrt3}
       \frac{(v+\delta)^{-1/2}-v^{-1/2}}{\delta}.
\]
The integral identity
\[
 \Gamma_\delta(u,v)
   =-\frac{\sqrt u}{3\sqrt3}\int_0^1(v+t\delta)^{-3/2}\dd t
\]
gives uniformly bounded $C^{20}$ norm on a slightly larger positive
rectangle, independently of small $\delta$. It also shows, for every
fixed $k$, that this is a $C^k$ perturbation, of size $O(\delta)$, of
\[
 \Gamma_0(u,v)=-\frac1{3\sqrt3}u^{1/2}v^{-3/2}.
\]
The exact $m,m$ derivative never vanishes: the phase is $\sqrt m$ times
a strictly negative function of $l$. The model Hessian determinant is
nonzero, hence the exact normalized determinant is uniformly nonzero
for sufficiently small $\delta$. We may therefore use the fixed-order
$\Rfine(M,L,P)$ of Lemma~\ref{lem:sharpB}.

Solving the two gradient equations of the model monomial gives a Legendre
phase
\[
 C_1|r|^{-1/4}|s|^{3/4},\qquad C_1\ne0,
\]
on its fixed opposite-sign frequency quadrant. On normalized compact
inverse environments, $|r|,|s|$ stay bounded away from zero. The pure third
derivative in $s$ is
\[
 \frac{15}{64}C_1\operatorname{sgn}(s)
                      |r|^{-1/4}|s|^{-9/4},
\]
which is uniformly nonzero. The common inverse argument of
Lemma~\ref{lem:transfer} transfers this to the exact phase.

If both dual lengths $P/M,P/L$ are at least one, use
\eqref{eq:third-test} in the second dual variable and count the first.
Exactly as in the proof of \eqref{eq:terminal}, this gives
$\sqrt L P^{2/3}+P^{5/6}$. Including the complete finer remainder yields
\begin{equation}\label{eq:B2}
 B_2(h)\ll x^\eta\left(
       \sqrt L P^{2/3}+P^{5/6}+\sqrt P+\frac{ML}{\sqrt P}+L\right).
\end{equation}
Here $ML/P$ was absorbed by $ML/\sqrt P$, since $P=hX\ge1$.

We verify \eqref{eq:B2} also when a dual side is short. If $P<L$, the
complete trivial dual main term
\[
 P+M+L+ML/P
\]
and $\Rfine(M,L,P)$ are bounded, up to a logarithm, by
$L+ML/\sqrt P$. Indeed $P\le L$, $L\ge\sqrt P$, and
$M\le ML/\sqrt P$.

The only remaining short-dual case is $L\le P<M$. Since $P=hX\ge X$,
it implies $X^2<N$. Consequently
\begin{equation}\label{eq:B-short-margin}
 \frac{M}{L^2}=\frac{NX}{F^2}
 \ll N^{5/6}x^{-2/3}
 \le x^{-287/2418},\qquad
 \frac{5\rho-4}{6}=-\frac{287}{2418}.
\end{equation}
For sufficiently large $x$, we have $P<M<L^2$, so again
$L\ge\sqrt P$. Both $P$ and $M$ are bounded by $ML/\sqrt P$, and the
same absorption applies. This argument uses the actual global range
\eqref{eq:caseB-range}; it is not asserted for arbitrary $M,L$.
These cases exhaust all short duals, including their equality boundaries,
and prove \eqref{eq:B2} for every nonempty correlation.

\subsection{Summation, caps and integer rounding}
The contribution of $ML/\sqrt{hX}$ to \eqref{eq:caseB-A}, divided by its
diagonal $M^2F/q$, is $O(\sqrt{q/X})$, hence is absorbed by
\eqref{eq:caseB-cap}. Also $\sqrt P\le P^{5/6}$. The remaining variable
terms for $q\ge2$ are
\begin{equation}\label{eq:B-balance}
 \frac aq+bq^{2/3}+c_1q^{5/6},\qquad
 a=\frac{N^2F}{X^2},\quad b=NF^{1/2}X^{1/6},\quad c_1=NX^{5/6},
\end{equation}
together with $NF/X+N^2/F+M^2$. For $q=1$, the corresponding estimate is
only $a+N^2/F+M^2$, with no correlation terms inserted.

Choose
\begin{equation}\label{eq:B-q}
 q_b=(a/b)^{3/5},\quad q_c=(a/c_1)^{6/11},\quad
 K=\frac{\min(X,L)}{C\log x},\quad
 q=\max\{1,\lfloor\min(q_b,q_c,K)\rfloor\}.
\end{equation}
If the minimum is at least two, $q$ is between half that minimum and the
minimum. Hence
\begin{align*}
 a/q&\ll a^{2/5}b^{3/5}+a^{5/11}c_1^{6/11}+a/K,\\
 bq^{2/3}&\le a^{2/5}b^{3/5},\qquad
 c_1q^{5/6}\le a^{5/11}c_1^{6/11}.
\end{align*}
The cap term satisfies
\[
 a/K\ll\log x\left(\frac{N^2}{X}+\frac{N^2F}{X^3}\right).
\]
If the minimum is less than two, $q=1$. The responsible term among
$a/q_b,a/q_c,a/K$ is already at least $a/2$, so the same retained bound
follows. This includes $K<1$ without requiring an impossible nonempty-shift
cap for $q=1$.

Substituting the two balances gives the seven terms
\begin{align}\label{eq:Bfinal}
 |S_1|^2\ll x^\eta\biggl(&\frac{N^2}{X}+\frac{N^2F}{X^3}
       +N^{7/5}F^{7/10}X^{-7/10}\notag\\
       &+N^{16/11}F^{5/11}X^{-5/11}
       +\frac{NF}{X}+\frac{N^2}{F}+M^2\biggr).
\end{align}
Write $N=x^u$, $X=x^v$ and $\Fzero=x^{(1+u)/3}$ as before. After taking
a square root and multiplying by $x^{1/3}N^{-2/3}$, the seven budget
exponents are
\begin{align}
 B_1&=\frac13+\frac u3-\frac v2,
 & B_2&=\frac12+\frac u2-\frac{3v}{2},\notag\\
 B_3&=\frac9{20}+\frac{3u}{20}-\frac{7v}{20},
 & B_4&=\frac9{22}+\frac{3u}{22}-\frac{5v}{22},
                 \label{eq:Bexponents}\\
 B_5&=\frac12-\frac v2,
 & B_6&=\frac{1+u}{6},\qquad
 B_7=\frac13+\frac u3-v.\notag
\end{align}
The last seven certificates in Appendix~\ref{app:affine} prove
$B_j<\theta$ on $u_0\le u\le\rho$, $9u/20\le v\le u$.
The smallest constant margin is $11/241800$. The bound
\eqref{eq:caseB-short-first} uses $B_6,B_7$, so it is covered as well.
The fixed relation $F=3\Fzero$ changes only constants. This proves
\eqref{eq:box-target} on the whole Case B range.

\section{Assembly of the proposed pointwise bound}
\label{sec:assembly}
We complete the argument for Theorem~\ref{thm:main}. A factor box with
$N\le x^{u_0}$ contains $O(N)$ triples, so its contribution after product
Abel is
\[
 O(x^{1/3}N^{1/3})\ll x^{1/3+u_0/3}=x^\theta.
\]
For every larger nonempty box, \eqref{eq:box-range} gives exactly
$u_0\le\log_xN\le\rho$ and $N^{1/3}\le X\le N$. The two cases
$X\le N^{9/20}$ and $X\ge N^{9/20}$ overlap at the boundary and exhaust
this interval. Sections~\ref{sec:caseA} and \ref{sec:caseB} prove
\eqref{eq:box-target} uniformly in the product cuts needed by Abel
summation. Summing the $O(\log^3x)$ factor boxes bounds the main oscillating
sum in \eqref{eq:voronoi}.

All original factor cutoffs, first-A subintervals, stationary domains,
rectangular cuts and shifted intersections lie in the uniform domain
families already treated. The estimates are for each original real $x$;
no averaging in $x$ or almost-everywhere assertion has been used. The
intermediate endpoint and isolated-fiber costs are those explicitly
retained in the B errors. In particular, the cutoff at $Y$ is part of
the product-domain estimate, not a discarded summand.

Choose one auxiliary $\eta>0$ sufficiently small compared with the desired
$\eps$. There are finitely many sums and Cauchy or Abel operations, and
only fixed logarithmic powers. Thus the total loss has the form
$x^{C\eta}(\log x)^C$ for a fixed $C$. Choose $\eta<\eps/(2C)$ and absorb
the logarithm into $x^{\eps/2}$. The positive margins establishing chart
admissibility are independent of this choice. Since
$Y=\lfloor x^\rho\rfloor\ge x^\rho/2$ for large $x$, the remaining tail is
\[
 x^{2/3+\eta}Y^{-1/3}
 \ll x^{2/3-\rho/3+\eta}=x^{\theta+\eta}.
\]
Together with \eqref{eq:endpoint-identities}, this is the proposed bound
$\Delta_3(x)\ll_\eps x^{541/1209+\eps}$. Finite smaller $x$ are covered by
enlarging the implied constant.\hfill$\square$

\appendix
\section{Exact model derivative certificates}
\label{app:model}
This appendix supplies the algebraic details used in
Lemma~\ref{lem:model}. All derivatives of Legendre phases are evaluated
at the successive points corresponding to the original $(u,v)$.
For a polynomial $P(z)=\sum_{j=0}^d c_jz^j$, write
\[
 P^{\rm hom}(a,b)=\sum_{j=0}^d c_j a^{d-j}b^j.
\]

\subsection{The quartic and degree-nine polynomials}
Define
\begin{align*}
 P_4(z)&=7+70z+252z^2+432z^3+324z^4,\\
 Q_4(z)&=1-8z-72z^2-108z^3-108z^4,\\
 R_9(z)&=-1-17z-196z^2-920z^3+468z^4+19332z^5\\
       &\quad+75600z^6+124416z^7+93312z^8+46656z^9,\\
 S_9(z)&=7+168z+1764z^2+11010z^3+45711z^4+131814z^5\\
       &\quad+265896z^6+363528z^7+306180z^8+122472z^9.
\end{align*}
At $u=v=1$, put
\[
 D_{ab}=a^2+4ab+6b^2,\qquad E_{ab}=a^2+6ab+18b^2.
\]
The exact pure derivatives are
\begin{align}
 K_{0,11}&=-\frac{16P_4^{\rm hom}(a,b)}{9D_{ab}^3},
 &K_{0,22}&=\frac{2Q_4^{\rm hom}(a,b)}{9D_{ab}^3},
                    \label{eq:Kquartics}\\
 \Psi_{0,111}&=\frac{27R_9^{\rm hom}(a,b)}{1600E_{ab}^3},
 &\Psi_{0,222}&=\frac{27S_9^{\rm hom}(a,b)}{25E_{ab}^3}.
                    \label{eq:Psi-at-one}
\end{align}
For general positive $u,v$, let $t=u/v$ and
$E_0=a^2+6abt+18b^2t^2$. The same pure third derivatives are
\begin{align}
 \Psi_{0,111}
   &=\frac{27a^9 R_9(bt/a)}{1600v^{13/2}t^{7/2}E_0^3},\notag\\
 \Psi_{0,222}
   &=\frac{27a^9 S_9(bt/a)}{25v^{13/2}\sqrt t\,E_0^3}.
                  \label{eq:Psi-general}
\end{align}
At $a=0$ these expressions mean their homogeneous polynomial extensions,
not division by $a$.

\subsection{A direct differentiation procedure}
Let $I=(\Hess G_0)^{-1}$. Since $\Phi_{0,p}=-u$, put
\begin{align}
 r=K_{0,p}=-I_{11}
   &=\frac{8u^{3/2}v^2(av+3bu)}{3D},\notag\\
 s=K_{0,q}=-I_{12}
   &=\frac{2\sqrt u\,v^3(av+6bu)}{3D}.
                  \label{eq:r-s-inverse}
\end{align}
The homogeneity degree of $G_0$ in $(u,v)$ is $-1/2$. Consequently
$\Phi_0$ has degree $1/3$ in $(p,q)$, and $K_0$ has degree $-2/3$.
Euler's identity therefore gives
\begin{equation}\label{eq:Psi-value}
 \Psi_0=K_0-pr-qs=\frac53 K_0=-\frac53u.
\end{equation}
Use the inverse Jacobian of $(u,v)\mapsto(r,s)$ in
\eqref{eq:r-s-inverse} to differentiate \eqref{eq:Psi-value} three times
in $r$ or $s$. This gives \eqref{eq:Psi-general}. In particular, this
procedure does not assume $r$ or $s$ is nonzero.

All expressions can be checked using rational functions by setting
$u=w^2$, so $\partial_u=(2w)^{-1}\partial_w$. Explicitly define
\[
 D_p=I_{11}\partial_u+I_{21}\partial_v,\qquad
 D_q=I_{12}\partial_u+I_{22}\partial_v.
\]
Then $r=D_p(-u)$, $s=D_q(-u)$ and
$\Hess K_0$ has entries $D_iD_j(-u)$. If
$J=\partial(r,s)/\partial(w,v)$, the columns of $J^{-1}$ give the
coefficients of $\partial_r,\partial_s$ as operators in
$\partial_w,\partial_v$. This finite differentiation procedure verifies
\eqref{eq:model-dets}, \eqref{eq:Kquartics} and
\eqref{eq:Psi-general} without explicitly solving a high-degree inverse.

There is also a tensor check at $(1,1)$. The jet is
\[
 \partial_u^j\partial_v^kG_0(1,1)
   =a(1/2)_{\underline j}(-1)_{\underline k}
       +b(3/2)_{\underline j}(-2)_{\underline k},
\]
where the underline denotes a falling factorial. Write $v_i=Ie_i$ and
let the vector $B_{ij}$ have components $G_0'''(e_r,v_i,v_j)$. Differentiating
$\Hess\Phi_0=-I$ gives
\begin{align*}
 \Phi_{0,ijk}&=G_0'''(v_i,v_j,v_k),\\
 \Phi_{0,ijkl}&=G_0''''(v_i,v_j,v_k,v_l)
       -B_{ij}^{\mathsf T}IB_{kl}-B_{ik}^{\mathsf T}IB_{jl}
       -B_{il}^{\mathsf T}IB_{jk}.
\end{align*}
Thus if $L_{ij}=\Phi_{0,1ij}$ and $J_0=L^{-1}$, then
\[
 \Psi_{0,rrr}=K_0'''(J_0e_r,J_0e_r,J_0e_r),
 \qquad K_{0,ijk}=\Phi_{0,1ijk}.
\]
Substitution of the displayed jet gives exactly
\eqref{eq:Kquartics} and \eqref{eq:Psi-at-one}.

\subsection{Noncoincidence and the missing projective direction}
For a list $[c_0,\ldots,c_d]$, let $[c_0,\ldots,c_d](z)$ mean
$\sum_{j=0}^d c_jz^j$. The following two polynomial identities hold in
$\mathbb F_{101}[z]$:
\begin{align}
 &[81,68,77,20](z)P_4(z)
       +[40,32,49,60](z)Q_4(z)=1,\label{eq:bezout4}\\
 &[8,56,98,2,87,80,95,42,3](z)R_9(z)\notag\\
 &\hspace{22mm}+[59,11,89,19,55,63,66,100,71](z)S_9(z)=1.
                         \label{eq:bezout9}
\end{align}
They are checked by polynomial multiplication and coefficient reduction.
The leading coefficients of $P_4,Q_4,R_9,S_9$ modulo $101$ are
$21,94,95,60$, respectively. Thus the degrees survive reduction and the
two pairs are coprime over $\Q$ (equivalently, their integer resultants
are nonzero; their residues are $97$ and $51$). Here and below $\Q$
denotes the field of rational numbers.

At $a=0$, $b\ne0$, the unit-point values are
\begin{equation}\label{eq:a-zero}
 K_{0,11}=-\frac8{3b^2},\quad K_{0,22}=-\frac1{9b^2},\quad
 \Psi_{0,111}=\frac{27b^3}{200},\quad
 \Psi_{0,222}=\frac{567b^3}{25}.
\end{equation}
This checks the projective direction omitted by setting $z=b/a$.

Finally, checking at $(1,1)$ loses no points. At any $(u_*,v_*)>0$, rescale
the analytic variables and replace the parameter pair by
\[
 (a',b')=\left(a\sqrt{u_*}/v_*,\ b u_*^{3/2}/v_*^2\right).
\]
The pure second derivatives of $K_0$ acquire the positive scaling factors
$u_*^3$ and $u_*v_*^2$, and the pure third derivatives of $\Psi_0$ acquire
$u_*^{-5}$ and $u_*^{-2}v_*^{-3}$. The transformed parameter pair stays
in a compact set separated from zero. Since $D_{ab}$ and $E_{ab}$ are
positive away from $(0,0)$, the polynomial noncoincidence, the explicit
axis values and compactness prove all uniform lower bounds asserted in
Lemma~\ref{lem:model}.

\subsection{Fourth and fifth derivatives in the same coordinate}
\label{app:four-five}
In addition to the degree-nine certificates above, we require the
same-coordinate condition \eqref{eq:pure-four-five-model}. Continuing
the inverse-Jacobian differentiation in \eqref{eq:r-s-inverse} in the
second terminal coordinate gives, at $u=v=1$,
\begin{align}
 \Psi_{0,2222}
   &=-\frac{81Q_{14}^{\rm hom}(a,b)}{250E_{ab}^5},\label{eq:Psi-four}\\
 \Psi_{0,22222}
   &=\frac{243Q_{19}^{\rm hom}(a,b)}{1250E_{ab}^7}.\label{eq:Psi-five}
\end{align}
The following table gives all coefficients in
$Q_d(z)=\sum_{i=0}^d c_i z^i$. The homogeneous convention is that
at the beginning of Appendix~\ref{app:model}.
\begin{center}
\small
\begin{longtable}{r rr}
\toprule
$i$ & $[z^i]Q_{14}$ & $[z^i]Q_{19}$\\
\midrule
\endfirsthead
\toprule
$i$ & $[z^i]Q_{14}$ & $[z^i]Q_{19}$\\
\midrule
\endhead
\bottomrule
\endfoot
0 & 119 & 1309\\
1 & 4522 & 68068\\
2 & 77826 & 1625778\\
3 & 824568 & 24203784\\
4 & 6106629 & 255656181\\
5 & 33762834 & 2057394312\\
6 & 144505512 & 13172633964\\
7 & 487756296 & 68944652064\\
8 & 1305974988 & 299918873502\\
9 & 2759927904 & 1094360860548\\
10 & 4528553832 & 3361429410120\\
11 & 5592351456 & 8683008133680\\
12 & 4917915648 & 18765796869360\\
13 & 2764437984 & 33609519267840\\
14 & 753937632 & 49138093018080\\
15 & 0 & 57316640278464\\
16 & 0 & 51469491145248\\
17 & 0 & 33532209482688\\
18 & 0 & 14167995980544\\
19 & 0 & 2931309513216\\
\end{longtable}
\end{center}

For a polynomial specified by its ascending coefficient list, there is
the identity
\begin{equation}\label{eq:bezout-four-five}
 A_{18}(z)Q_{14}(z)+B_{13}(z)Q_{19}(z)\equiv1\pmod{101}.
\end{equation}
To specify the certificate without a long unbroken coefficient list, write
\[
 A_{18}(z)=A^{(0)}(z)+z^{10}A^{(1)}(z),\qquad
 B_{13}(z)=B^{(0)}(z)+z^7B^{(1)}(z),
\]
where the ascending lists are
\begin{align*}
 A^{(0)}&=[45,18,2,88,53,55,93,1,71,67],\\
 A^{(1)}&=[100,76,17,21,72,87,0,57,6],\\
 B^{(0)}&=[76,87,44,85,8,90,22],\\
 B^{(1)}&=[11,90,96,45,74,100,12].
\end{align*}
All coefficients are interpreted in $\mathbb F_{101}$. Direct
multiplication verifies \eqref{eq:bezout-four-five}. The leading
coefficients of $Q_{14}$ and $Q_{19}$ reduce to $3$ and $49$, so their
degrees survive reduction. Their resultant is $29$ modulo $101$.
Consequently the two integer polynomials are coprime over $\mathbb Q$
and have no common finite complex zero.

At $a=0$, the omitted projective direction is
\begin{equation}\label{eq:four-five-axis}
 \Psi_{0,2222}=-\frac{32319}{250}b^4,\qquad
 \Psi_{0,22222}=\frac{581742}{625}b^5,
\end{equation}
which are both nonzero for $b\ne0$.
Under the positive-coordinate rescaling described above, the fourth
and fifth derivatives acquire respectively the positive factors
$u_*^{-3}v_*^{-4}$ and $u_*^{-4}v_*^{-5}$.
The transformed parameters remain in a compact set separated from zero,
and $E_{ab}>0$ there. Compactness gives the uniform lower bound
\eqref{eq:pure-four-five-model}. The $C^5$ comparison in
Lemma~\ref{lem:transfer} transfers it to the exact phase. Thus zeros
of its fourth derivative are treated by Lemma~\ref{lem:sublevel-fourth}
without replacing the phase in an exponential sum by its model.

\section{Exact affine certificates for all parameter ranges}
\label{app:affine}
All identities in this appendix are over the rational numbers. They can be
verified by comparing coefficients of $1,u,v,h$. Every right side is a
nonnegative rational combination of generators that are nonnegative on
the stated region. No numerical vertex search is needed to use the
certificates.

\subsection{Deficits and common generators}
Retain
\[
 u_0=\frac{138}{403},\quad \rho=\frac{265}{403},\quad
 f=\frac{1+u}{3},\quad
 \hthreshold=\frac{2u-f-v}{3},\quad
 \hcap=\frac{2(u-u_0)}3,
\]
and $c,\mathfrak p_0,\mathfrak p_b$ from \eqref{eq:c-p0-pb}.
Write $\mathfrak p=\log_x\Pi$, $\mathfrak a=\log_x\alpha$ for the
nominal, not rounded, scales. In the axis and main full regions use
$(a_j,b_j,d_j)=(a_j^{\mathrm L},b_j^{\mathrm L},d_j^{\mathrm L})$
from \eqref{eq:local-powers}; in the boundary full region use the
conservative row $\star=\mathrm S$.

The propagated deficits for \eqref{eq:A2} followed by \eqref{eq:A1} are
\begin{align}
 \mathcal D&=u-\frac{f+3h+v-\mathfrak p}{2},\label{eq:deficit-D}\\
 \mathcal C_j&=u-\frac{f+h-v}{2}
 -\frac{h+v+a_j(f+h-v+\mathfrak a)+b_jh+d_jv}{2}.
 \label{eq:deficit-J}
\end{align}
In particular, $\mathcal C_j$ is a deficit for the contribution to
$S_2$, not merely for $\mathcal J_j$. Define
\[
 \mathcal E_{\rm end}=u-h-v,\qquad
 \mathcal E_{\rm curv}=\frac{f-h-v}{2},\qquad
 \Qone=\frac{\mathfrak p+h-v}{2}.
\]
The common generators are
\begin{align*}
 g_0&=1,&g_1&=u-u_0,&g_2&=\rho-u,\\
 g_3&=v-u/3,&g_4&=9u/20-v,&g_5&=h-\hthreshold,\\
 g_6&=\hcap-h.
\end{align*}
They are nonnegative on \eqref{eq:largeH-region}. The extra generators
and the selected exponents are specified in each subregion below.

\subsection{Case A: the axis region}
Here $c\le0$ and the long-direction powers are used:
\[
 g_7=-c,\qquad
 \mathfrak p=\frac{-h+92v-26f}{76},\qquad
 \mathfrak a=\mathfrak p-v.
\]
The 14 certificates are
\begingroup
\small
\begin{align}
% CERT|axis|D|329/24180*g0 + 6/35*g3 + 67/56*g6 + 11/2660*g7
\mathcal D &= \frac{329}{24180}g_{0}+\frac{6}{35}g_{3}+\frac{67}{56}g_{6}+\frac{11}{2660}g_{7}\label{cert:axis:1}\\
% CERT|axis|C1|329/24180*g0 + 6/35*g3 + 67/56*g6 + 11/2660*g7
\mathcal C_1 &= \frac{329}{24180}g_{0}+\frac{6}{35}g_{3}+\frac{67}{56}g_{6}+\frac{11}{2660}g_{7}\label{cert:axis:2}\\
% CERT|axis|C2|941/64480*g0 + 11/140*g3 + 521/448*g6 + 87/21280*g7
\mathcal C_2 &= \frac{941}{64480}g_{0}+\frac{11}{140}g_{3}+\frac{521}{448}g_{6}+\frac{87}{21280}g_{7}\label{cert:axis:3}\\
% CERT|axis|C3|5675/251472*g0 + 5/182*g4 + 3287/2912*g6 + 87/55328*g7
\mathcal C_3 &= \frac{5675}{251472}g_{0}+\frac{5}{182}g_{4}+\frac{3287}{2912}g_{6}+\frac{87}{55328}g_{7}\label{cert:axis:4}\\
% CERT|axis|C4|4817/125736*g0 + 5/91*g4 + 1649/1456*g6 + 87/27664*g7
\mathcal C_4 &= \frac{4817}{125736}g_{0}+\frac{5}{91}g_{4}+\frac{1649}{1456}g_{6}+\frac{87}{27664}g_{7}\label{cert:axis:5}\\
% CERT|axis|C5|17/3720*g0 + 29/70*g3 + 143/112*g6 + 13/10640*g7
\mathcal C_5 &= \frac{17}{3720}g_{0}+\frac{29}{70}g_{3}+\frac{143}{112}g_{6}+\frac{13}{10640}g_{7}\label{cert:axis:6}\\
% CERT|axis|C6|15403/502944*g0 + 25/52*g4 + 757/832*g6 + 1/832*g7
\mathcal C_6 &= \frac{15403}{502944}g_{0}+\frac{25}{52}g_{4}+\frac{757}{832}g_{6}+\frac{1}{832}g_{7}\label{cert:axis:7}\\
% CERT|axis|C7|11899/335296*g0 + 445/728*g4 + 9897/11648*g6 + 1191/221312*g7
\mathcal C_7 &= \frac{11899}{335296}g_{0}+\frac{445}{728}g_{4}+\frac{9897}{11648}g_{6}+\frac{1191}{221312}g_{7}\label{cert:axis:8}\\
% CERT|axis|p-1/1000|826529/22971000*g0 + 15/19*g3 + 24/19*g5 + 97/76*g6
\mathfrak p-\frac1{1000} &= \frac{826529}{22971000}g_{0}+\frac{15}{19}g_{3}+\frac{24}{19}g_{5}+\frac{97}{76}g_{6}\label{cert:axis:9}\\
% CERT|axis|-alpha-1/1000|132393863/1079637000*g0 + 220/893*g4 + 96/893*g5 + 337/3572*g6
-\mathfrak a-\frac1{1000} &= \frac{132393863}{1079637000}g_{0}+\frac{220}{893}g_{4}+\frac{96}{893}g_{5}+\frac{337}{3572}g_{6}\label{cert:axis:10}\\
% CERT|axis|f+h-u-1/1000|353789/5211000*g0 + 80/193*g4 + 2912/1737*g5 + 47/5211*g7
f+h-u-\frac1{1000} &= \frac{353789}{5211000}g_{0}+\frac{80}{193}g_{4}+\frac{2912}{1737}g_{5}+\frac{47}{5211}g_{7}\label{cert:axis:11}\\
% CERT|axis|v-h-1/1000|4007/93000*g0 + 39/35*g3 + 13/28*g6 + 1/140*g7
v-h-\frac1{1000} &= \frac{4007}{93000}g_{0}+\frac{39}{35}g_{3}+\frac{13}{28}g_{6}+\frac{1}{140}g_{7}\label{cert:axis:12}\\
% CERT|axis|Berr1|1349/8112*g0 + 25/26*g4 + 341/416*g6 + 1/416*g7
\mathcal E_{\rm end} &= \frac{1349}{8112}g_{0}+\frac{25}{26}g_{4}+\frac{341}{416}g_{6}+\frac{1}{416}g_{7}\label{cert:axis:13}\\
% CERT|axis|Berr2|359/5211*g0 + 80/193*g4 + 613/3474*g5 + 47/5211*g7
\mathcal E_{\rm curv} &= \frac{359}{5211}g_{0}+\frac{80}{193}g_{4}+\frac{613}{3474}g_{5}+\frac{47}{5211}g_{7}\label{cert:axis:14}
\end{align}
\endgroup


\subsection{Case A: the main full region}
Here $c\ge0$, $\mathfrak p_0\ge\mathfrak p_b$, and the long-direction
powers are used:
\[
 g_7=c,\quad g_8=\mathfrak p_0-\mathfrak p_b,\quad
 \mathfrak p=\mathfrak p_0,\quad
 \mathfrak a=\frac{\mathfrak p-h-v}{2}.
\]
The 15 certificates are
\begingroup
\small
\begin{align}
% CERT|main|D|7/54*g2 + 8/63*g3 + 59/42*g6
\mathcal D &= \frac{7}{54}g_{2}+\frac{8}{63}g_{3}+\frac{59}{42}g_{6}\label{cert:main:1}\\
% CERT|main|C1|7/54*g2 + 8/63*g3 + 59/42*g6
\mathcal C_1 &= \frac{7}{54}g_{2}+\frac{8}{63}g_{3}+\frac{59}{42}g_{6}\label{cert:main:2}\\
% CERT|main|C2|5/19344*g0 + 59/432*g2 + 2/63*g3 + 929/672*g6
\mathcal C_2 &= \frac{5}{19344}g_{0}+\frac{59}{432}g_{2}+\frac{2}{63}g_{3}+\frac{929}{672}g_{6}\label{cert:main:3}\\
% CERT|main|C3|35849/1695018*g0 + 15/4907*g4 + 42979/39256*g6 + 7/1402*g8
\mathcal C_3 &= \frac{35849}{1695018}g_{0}+\frac{15}{4907}g_{4}+\frac{42979}{39256}g_{6}+\frac{7}{1402}g_{8}\label{cert:main:4}\\
% CERT|main|C4|10333/437658*g0 + 43/543*g2 + 827/724*g6 + 3/181*g8
\mathcal C_4 &= \frac{10333}{437658}g_{0}+\frac{43}{543}g_{2}+\frac{827}{724}g_{6}+\frac{3}{181}g_{8}\label{cert:main:5}\\
% CERT|main|C5|17/3720*g0 + 29/70*g3 + 143/112*g6 + 1/720*g7
\mathcal C_5 &= \frac{17}{3720}g_{0}+\frac{29}{70}g_{3}+\frac{143}{112}g_{6}+\frac{1}{720}g_{7}\label{cert:main:6}\\
% CERT|main|C6|41099/1695018*g0 + 260/701*g4 + 2123/2804*g6 + 63/2804*g8
\mathcal C_6 &= \frac{41099}{1695018}g_{0}+\frac{260}{701}g_{4}+\frac{2123}{2804}g_{6}+\frac{63}{2804}g_{8}\label{cert:main:7}\\
% CERT|main|C7|1/8*g8
\mathcal C_7 &= \frac{1}{8}g_{8}\label{cert:main:8}\\
% CERT|main|p-1/1000|4007/93000*g0 + 39/35*g3 + 13/28*g6 + 11/1260*g7
\mathfrak p-\frac1{1000} &= \frac{4007}{93000}g_{0}+\frac{39}{35}g_{3}+\frac{13}{28}g_{6}+\frac{11}{1260}g_{7}\label{cert:main:9}\\
% CERT|main|-alpha-1/1000|1963783/15717000*g0 + 20/91*g4 + 11/364*g6 + 19/3276*g7
-\mathfrak a-\frac1{1000} &= \frac{1963783}{15717000}g_{0}+\frac{20}{91}g_{4}+\frac{11}{364}g_{6}+\frac{19}{3276}g_{7}\label{cert:main:10}\\
% CERT|main|f+h-u-1/1000|179623/3627000*g0 + 47/180*g2 + 1/3*g4 + 1*g5
f+h-u-\frac1{1000} &= \frac{179623}{3627000}g_{0}+\frac{47}{180}g_{2}+\frac{1}{3}g_{4}+g_{5}\label{cert:main:11}\\
% CERT|main|v-h-1/1000|9791/1209000*g0 + 1/3*g2 + 1*g3 + 1*g6
v-h-\frac1{1000} &= \frac{9791}{1209000}g_{0}+\frac{1}{3}g_{2}+g_{3}+g_{6}\label{cert:main:12}\\
% CERT|main|Berr1|130126/847509*g0 + 520/701*g4 + 721/1402*g6 + 63/1402*g8
\mathcal E_{\rm end} &= \frac{130126}{847509}g_{0}+\frac{520}{701}g_{4}+\frac{721}{1402}g_{6}+\frac{63}{1402}g_{8}\label{cert:main:13}\\
% CERT|main|Berr2|166567/6586632*g0 + 22193/81720*g2 + 319/1362*g4 + 21/454*g8
\mathcal E_{\rm curv} &= \frac{166567}{6586632}g_{0}+\frac{22193}{81720}g_{2}+\frac{319}{1362}g_{4}+\frac{21}{454}g_{8}\label{cert:main:14}\\
% CERT|main|q1|1/126*g7
\Qone &= \frac{1}{126}g_{7}\label{cert:main:15}
\end{align}
\endgroup


The first two equalities are the active balances. In particular,
\[
 \mathcal D=\mathcal C_1
 =\frac7{54}(\rho-u)+\frac8{63}(v-u/3)
     +\frac{59}{42}(\hcap-h).
\]
The second correlation has the strict margin
\[
 \mathcal C_2=\frac5{19344}+\frac{59}{432}(\rho-u)
 +\frac2{63}(v-u/3)+\frac{929}{672}(\hcap-h).
\]

\subsection{Case A: the boundary full region}
Here $c\ge0$, $\mathfrak p_b\ge\mathfrak p_0$, and the conservative
powers from \eqref{eq:L5safe} must be used:
\[
 g_7=c,\quad g_8=\mathfrak p_b-\mathfrak p_0,\quad
 \mathfrak p=\mathfrak p_b,\quad
 \mathfrak a=\frac{\mathfrak p-h-v}{2}.
\]
The last correlation deficit is zero, as expected from $\Pi=\Pi_b$.
The 15 certificates are
\begingroup
\small
\begin{align}
% CERT|boundary|D|640/72943*g0 + 187/1086*g2 + 595/362*g6 + 189/362*g8
\mathcal D &= \frac{640}{72943}g_{0}+\frac{187}{1086}g_{2}+\frac{595}{362}g_{6}+\frac{189}{362}g_{8}\label{cert:boundary:1}\\
% CERT|boundary|C1|149/241800*g0 + 1331/3000*g2 + 33/50*g4 + 277/100*g6
\mathcal C_1 &= \frac{149}{241800}g_{0}+\frac{1331}{3000}g_{2}+\frac{33}{50}g_{4}+\frac{277}{100}g_{6}\label{cert:boundary:2}\\
% CERT|boundary|C2|15329/13540800*g0 + 72251/168000*g2 + 1593/2800*g4 + 3673/1400*g6
\mathcal C_2 &= \frac{15329}{13540800}g_{0}+\frac{72251}{168000}g_{2}+\frac{1593}{2800}g_{4}+\frac{3673}{1400}g_{6}\label{cert:boundary:3}\\
% CERT|boundary|C3|23/1488*g0 + 27/80*g2 + 3/4*g4 + 5/2*g6
\mathcal C_3 &= \frac{23}{1488}g_{0}+\frac{27}{80}g_{2}+\frac{3}{4}g_{4}+\frac{5}{2}g_{6}\label{cert:boundary:4}\\
% CERT|boundary|C4|1771/80600*g0 + 847/3000*g2 + 21/50*g4 + 199/100*g6
\mathcal C_4 &= \frac{1771}{80600}g_{0}+\frac{847}{3000}g_{2}+\frac{21}{50}g_{4}+\frac{199}{100}g_{6}\label{cert:boundary:5}\\
% CERT|boundary|C5|41/1209*g0 + 8/15*g2 + 1*g4 + 31/8*g6
\mathcal C_5 &= \frac{41}{1209}g_{0}+\frac{8}{15}g_{2}+g_{4}+\frac{31}{8}g_{6}\label{cert:boundary:6}\\
% CERT|boundary|C6|75/3224*g0 + 7/120*g2 + 1/2*g4 + 1*g6
\mathcal C_6 &= \frac{75}{3224}g_{0}+\frac{7}{120}g_{2}+\frac{1}{2}g_{4}+g_{6}\label{cert:boundary:7}\\
% CERT|boundary|C7|0*g0
\mathcal C_7 &= 0g_{0}\label{cert:boundary:8}\\
% CERT|boundary|p-1/1000|12582057/72943000*g0 + 7/543*g2 + 785/362*g6 + 441/362*g8
\mathfrak p-\frac1{1000} &= \frac{12582057}{72943000}g_{0}+\frac{7}{543}g_{2}+\frac{785}{362}g_{6}+\frac{441}{362}g_{8}\label{cert:boundary:9}\\
% CERT|boundary|-alpha-1/1000|55847/403000*g0 + 61/60*g2 + 3*g4 + 5*g6
-\mathfrak a-\frac1{1000} &= \frac{55847}{403000}g_{0}+\frac{61}{60}g_{2}+3g_{4}+5g_{6}\label{cert:boundary:10}\\
% CERT|boundary|f+h-u-1/1000|82328171/823329000*g0 + 4907/20430*g2 + 362/681*g4 + 21/227*g8
f+h-u-\frac1{1000} &= \frac{82328171}{823329000}g_{0}+\frac{4907}{20430}g_{2}+\frac{362}{681}g_{4}+\frac{21}{227}g_{8}\label{cert:boundary:11}\\
% CERT|boundary|v-h-1/1000|16892171/218829000*g0 + 121/181*g2 + 1043/362*g6 + 63/362*g8
v-h-\frac1{1000} &= \frac{16892171}{218829000}g_{0}+\frac{121}{181}g_{2}+\frac{1043}{362}g_{6}+\frac{63}{362}g_{8}\label{cert:boundary:12}\\
% CERT|boundary|Berr1|733/4836*g0 + 7/60*g2 + 1*g4 + 1*g6
\mathcal E_{\rm end} &= \frac{733}{4836}g_{0}+\frac{7}{60}g_{2}+g_{4}+g_{6}\label{cert:boundary:13}\\
% CERT|boundary|Berr2|75/3224*g0 + 47/120*g2 + 1/2*g4 + 1/2*g6
\mathcal E_{\rm curv} &= \frac{75}{3224}g_{0}+\frac{47}{120}g_{2}+\frac{1}{2}g_{4}+\frac{1}{2}g_{6}\label{cert:boundary:14}\\
% CERT|boundary|q1|11899/282503*g0 + 3560/4907*g4 + 9897/9814*g6 + 909/1402*g8
\Qone &= \frac{11899}{282503}g_{0}+\frac{3560}{4907}g_{4}+\frac{9897}{9814}g_{6}+\frac{909}{1402}g_{8}\label{cert:boundary:15}
\end{align}
\endgroup


\subsection{Case B}
On $u_0\le u\le\rho$, $9u/20\le v\le u$, take
\[
 b_0=1,\quad b_1=u-u_0,\quad b_2=\rho-u,\quad
 b_3=v-9u/20,\quad b_4=u-v.
\]
For the seven exponents in \eqref{eq:Bexponents}, the certificates are
\begingroup
\small
\begin{align}
% CERT|case_b|theta-B1|415/9672*b0 + 13/120*b2 + 1/2*b3
\theta-B_1 &= \frac{415}{9672}b_{0}+\frac{13}{120}b_{2}+\frac{1}{2}b_{3}\label{cert:case_b:1}\\
% CERT|case_b|theta-B2|179/24180*b0 + 7/40*b1 + 3/2*b3
\theta-B_2 &= \frac{179}{24180}b_{0}+\frac{7}{40}b_{1}+\frac{3}{2}b_{3}\label{cert:case_b:2}\\
% CERT|case_b|theta-B3|11/241800*b0 + 3/400*b1 + 7/20*b3
\theta-B_3 &= \frac{11}{241800}b_{0}+\frac{3}{400}b_{1}+\frac{7}{20}b_{3}\label{cert:case_b:3}\\
% CERT|case_b|theta-B4|1699/106392*b0 + 3/88*b2 + 5/22*b3
\theta-B_4 &= \frac{1699}{106392}b_{0}+\frac{3}{88}b_{2}+\frac{5}{22}b_{3}\label{cert:case_b:4}\\
% CERT|case_b|theta-B5|593/24180*b0 + 9/40*b1 + 1/2*b3
\theta-B_5 &= \frac{593}{24180}b_{0}+\frac{9}{40}b_{1}+\frac{1}{2}b_{3}\label{cert:case_b:5}\\
% CERT|case_b|theta-B6|69/403*b0 + 1/6*b2
\theta-B_6 &= \frac{69}{403}b_{0}+\frac{1}{6}b_{2}\label{cert:case_b:6}\\
% CERT|case_b|theta-B7|621/4030*b0 + 7/60*b1 + 1*b3
\theta-B_7 &= \frac{621}{4030}b_{0}+\frac{7}{60}b_{1}+b_{3}\label{cert:case_b:7}
\end{align}
\endgroup


There are $14+15+15+7=51$ identities. Each Case B constant term is
strictly positive, and the smallest is $11/241800$.
Definitions \eqref{eq:deficit-D}--\eqref{eq:deficit-J} connect all
certificates to the chosen local monomials and to both first-B errors.
Together with the domain, range and integer-rounding arguments in the
main text, they complete the parameter verification.


\begin{thebibliography}{9}
\bibitem{IZ}
A.~Ivi\'c and W.~Zhai,
\emph{Higher moments of the error term in the divisor problem},
arXiv:0904.2271v3 (2010), Lemma~3, pp.~6--8.
\url{https://arxiv.org/abs/0904.2271}.

\bibitem{Vand}
J.~Vandehey,
\emph{Error term improvements for van der Corput transforms},
arXiv:1205.0090v1 (2012), Theorem~1.1.
\url{https://arxiv.org/abs/1205.0090}.

\bibitem{HB}
D.~R.~Heath-Brown,
\emph{A new $k$-th derivative estimate for exponential sums via
Vinogradov's mean value},
arXiv:1601.04493v3 (2016), introduction, equation~(1).
\url{https://arxiv.org/abs/1601.04493}.

\bibitem{Basu}
S.~Basu,
\emph{Algorithmic Semi-Algebraic Geometry and Topology},
survey, Section~2, including Theorem~2.3 and the discussion of projection
and quantifier elimination.
\url{https://www.math.purdue.edu/~sbasu/ams_final_submission.pdf}.
\end{thebibliography}
\end{document}
